% !TeX root = linear-algebra.tex

\subsection{Differentiation}
\begin{definition}
For $A,B\in M_{m\times n}(\mathbb{R})$, the Frobenius inner product is
\[
    \innerproduct{A}{B}_F
    =
    \tr(A^{\top} B).
\]
The associated Frobenius norm is
$\norm{A}_F=\sqrt{\tr(A^{\top} A)}$.
\end{definition}

\begin{remark}\ \vspace{-2em}
\begin{enumerate}
    \item When $n=1$, the matrices $A,B$ are column vectors, and the Frobenius inner product reduces to the usual dot product.
    \item If $A=(A_{ij})$ and $B=(B_{ij})$, then
    \[
        \tr(A^{\top} B)
        =
        \sum_{j=1}^n\sum_{i=1}^m A_{ij}B_{ij}.
    \]
    In words, the Frobenius inner product is the ordinary dot product after writing the matrix entries as one long vector.
\end{enumerate}
\end{remark}


\begin{definition}[Derivative]
Let $f:M_{m\times n}(\mathbb{R})\to\mathbb{R}$ be differentiable. The derivative of $f$ at $X$ is the linear map
\[
    df_X:M_{m\times n}(\mathbb{R})\to\mathbb{R}
\]
characterized by
\[
    f(X+H)-f(X)=df_X(H)+r_X(H),
\]
Here $df_X(H)$ captures the terms which are linear in $H$.
The remainder $r_X(H)$ is of size $o(\norm{H}_F)$, i.e.
\[
    \frac{|r_X(H)|}{\norm{H}_F}\to 0
    \text{ as }H\to 0.
\]
\end{definition}

\begin{remark}
This is the matrix version of the definition of differentiability in multivariable calculus: near $X$, the function $f$ should be well approximated by a linear map, with relative error tending to zero.
In symbols, this means
\[
    \lim_{H\to 0}
    \frac{|f(X+H)-f(X)-df_X(H)|}{\norm{H}_F}
    =0.
\]
\end{remark}

\begin{example}
Consider $f:\mathbb{R}^2\to\mathbb{R}$ defined by $f(x_1,x_2)=x_1^2x_2$.
Writing $\mathbf{x}=(x_1,x_2)^{\top}$ and $\mathbf{h}=(h_1,h_2)^{\top}$, we have
\[
\begin{aligned}
    f(\mathbf{x}+\mathbf{h})-f(\mathbf{x})
    &=
    (x_1+h_1)^2(x_2+h_2)-x_1^2x_2 \\
    &=
    2x_1x_2h_1+x_1^2h_2
    +2x_1h_1h_2+x_2h_1^2+h_1^2h_2.
\end{aligned}
\]
The linear part in $\mathbf{h}$ is $2x_1x_2h_1+x_1^2h_2$, while the remaining terms are higher order. Hence
\[
    df_{\mathbf{x}}(\mathbf{h})
    =
    2x_1x_2h_1+x_1^2h_2
    =
    \left\langle
    \begin{pmatrix}
        2x_1x_2\\
        x_1^2
    \end{pmatrix},
    \begin{pmatrix}
        h_1\\
        h_2
    \end{pmatrix}
    \right\rangle.
\]
Therefore, $\nabla_{\mathbf{x}} f = (2x_1x_2, x_1^2)^{\top}$.
\end{example}

\begin{definition}
Let $R(H)$ be a scalar-valued function of a matrix variable
$H\in M_{m\times n}(\mathbb{R})$.

We write $R(H)=O(\norm{H}_F)$ as $H\to \mathbf O$ if there exist constants
$C>0$ and $\delta>0$ such that
\[
    |R(H)|
    \le
    C\norm{H}_F
    \qquad
    \text{whenever }0<\norm{H}_F<\delta.
\]
In words, $R(H)=O(\norm{H}_F)$ means that $R(H)$ is at most of first order in
the size of $H$.

We write $R(H)=o(\norm{H}_F)$ as $H\to \mathbf O$ if
\[
    \frac{|R(H)|}{\norm{H}_F}
    \to
    0
    \qquad
    \text{as }H\to \mathbf O.
\]
In words, $R(H)=o(\norm{H}_F)$ means that $R(H)$ is negligible compared with
the size of $H$.
\end{definition}

\begin{remark}
The same notation is often used for a one-dimensional perturbation $t\in
\mathbb{R}$. Since $t$ may be regarded as a $1\times1$ matrix, its Frobenius
norm is $\norm{t}_F=|t|$. Thus $R(t)=O(t)$ means $|R(t)|\le C|t|$ for all
sufficiently small nonzero $t$, and $R(t)=o(t)$ means
\[
    \frac{|R(t)|}{|t|}\to0
    \qquad
    \text{as }t\to0.
\]
This is the notation used when we restrict a matrix problem to a line such as
$X+tH$.
\end{remark}

\begin{example}
We have $\norm{H}_F^2=o(\norm{H}_F)$ because
\[
    \frac{\norm{H}_F^2}{\norm{H}_F}
    =
    \norm{H}_F
    \to
    0
    \qquad
    \text{as }H\to \mathbf O.
\]
On the other hand, $7\norm{H}_F=O(\norm{H}_F)$ because
$7\norm{H}_F\le 7\norm{H}_F$, but $7\norm{H}_F\neq o(\norm{H}_F)$ since
\[
    \frac{7\norm{H}_F}{\norm{H}_F}
    =
    7
    \not\to
    0.
\]
\end{example}

\begin{definition}[Gradient]
Let $f:M_{m\times n}(\mathbb{R})\to\mathbb{R}$ be differentiable. 
The unique matrix $G\in M_{m\times n}(\mathbb{R})$ such that
\[
    df_X(H)=\innerproduct{G}{H}_F=\tr(G^{\top} H)
    \quad\text{for all }H\in M_{m\times n}(\mathbb{R}).
\]

is called the gradient of $f$ at $X$, and is denoted by $\nabla_X f$.
\label{grad_def}
\end{definition}

Since $df_X$ is a linear functional on the inner product space $M_{m\times n}(\mathbb{R})$, the Riesz representation theorem gives a unique matrix $G\in M_{m\times n}(\mathbb{R})$ such that the definition~\eqref{grad_def} is valid.
Thus the gradient is characterized by
\[
    df_X(H)=\tr\bigl((\nabla_X f)^{\top} H\bigr).
\]


Informally, one may think of $dX$ as representing a small change $H$ in the matrix $X$. With this convention, the identity above is written as
\[
    df=\tr\bigl((\nabla_X f)^{\top} dX\bigr).
\]
The above ''equality'' gives the differential notation.

\begin{remark}\ \vspace{-2em}
\begin{enumerate}
    \item For $f:\mathbb{R}^n\to\mathbb{R}$, the gradient is the column vector collecting all partial derivatives:
    \[
        \nabla_{\mathbf{x}}f
        =
        \left(
            \frac{\partial f}{\partial x_1},\ldots,
            \frac{\partial f}{\partial x_n}
        \right)^{\top},
        \qquad
        df_{\mathbf{x}}(\mathbf{h})
        =
        (\nabla_{\mathbf{x}}f)^{\top}\mathbf{h}.
    \]
    Along a differentiable curve $\mathbf{x}(t)$, the chain rule gives
    \[
        \frac{d}{dt}f(\mathbf{x}(t))
        =
        df_{\mathbf{x}(t)}(\mathbf{x}'(t))
        =
        \sum_{i=1}^n
        \frac{\partial f}{\partial x_i}\frac{dx_i}{dt}.
    \]
    Writing $df=\frac{d}{dt}f\,dt$ and $dx_i=\frac{dx_i}{dt}\,dt$ yields
    \[
        df=(\nabla_{\mathbf{x}}f)^{\top}d\mathbf{x}
        =\sum_{i=1}^n\frac{\partial f}{\partial x_i}\,dx_i.
    \]
    Thus the differential formula is precisely the chain rule written without the parameter $t$.

    \item The derivative $df_X$ is a linear map. When $f$ is scalar-valued, it is a linear functional and the gradient is its Frobenius inner-product representative. Derivatives and differentials apply to scalar-, vector-, and matrix-valued functions, whereas a gradient is reserved for scalar-valued functions.
\end{enumerate}
\end{remark}

\begin{theorem}[Interior Extrema]
Let $U\subseteq M_{m\times n}(\mathbb{R})$ be open, and let $f:U\to\mathbb{R}$ be differentiable.
If $X_0\in U$ is a local maximum or a local minimum of $f$, then
\[
    \nabla_X f(X_0)=0.
\]
\end{theorem}

\begin{pfbox}
Fix any direction $H\in M_{m\times n}(\mathbb{R})$. Define a one-variable function
$\varphi(t)=f(X_0+tH)$ for $t$ near $0$.
Since $X_0$ is an interior local maximum or local minimum of $f$, the point $t=0$ is a local maximum or local minimum of $\varphi$. By Calculus 1, $\varphi'(0)=0$.

By the definition of the derivative,
\[
    \varphi'(0)=df_{X_0}(H).
\]
Thus $df_{X_0}(H)=0$ for every direction $H$.
Since
\[
    df_{X_0}(H)=\tr\bigl((\nabla_X f(X_0))^{\top}H\bigr),
\]
we have
\[
    \tr\bigl((\nabla_X f(X_0))^{\top}H\bigr)=0
    \qquad\text{for all }H.
\]
Taking $H=\nabla_X f(X_0)$ gives
\[
    \norm{\nabla_X f(X_0)}_F^2=0.
\]
Hence $\nabla_X f(X_0)=0$.
\end{pfbox}


Let $A=(A_{ij})$ be an $m\times n$ matrix whose entries are scalar-valued differentiable functions. We call $A$ a matrix-valued function.

\begin{definition}[Matrix Differential]
Let $A=(A_{ij})$ be an $m\times n$ matrix-valued function. Its differential is the $m\times n$ matrix
\[
    dA
    \overset{\mathrm{def}}{=}
    (dA_{ij})_{i,j}
    =
    \begin{pmatrix}
        dA_{11} & dA_{12} & \cdots & dA_{1n}\\
        dA_{21} & dA_{22} & \cdots & dA_{2n}\\
        \vdots & \vdots & \ddots & \vdots\\
        dA_{m1} & dA_{m2} & \cdots & dA_{mn}
    \end{pmatrix}.
\]
\end{definition}

\begin{example}
Let
\[
    A(x,y)
    =
    \begin{pmatrix}
        x^2y & e^x+y\\
        \sin x & x-y^2
    \end{pmatrix}.
\]
Then
\[
    dA
    =
    \begin{pmatrix}
        2xy\,dx+x^2\,dy & e^x\,dx+dy\\
        \cos x\,dx & dx-2y\,dy
    \end{pmatrix}.
\]
\end{example}


\begin{theorem}[Basic Differential Rules]
Let $A$ and $B$ be matrices of functions with compatible sizes, and let $\alpha\in\mathbb{R}$ be a constant scalar. Whenever the expressions are defined,
\begin{enumerate}
    \item $d(A+B)=dA+dB$;
    \item $d(\alpha A)=\alpha\,dA$;
    \item $d(AB)=(dA)B+A(dB)$;
    \item $d(A^{\top})=(dA)^{\top}$;
    \item $d\tr(A)=\tr(dA)$.
\end{enumerate}
\end{theorem}

\begin{pfbox}
The first two rules follow entry by entry from the linearity of ordinary differentials.
For the product rule, the $(i,j)$-entry of $AB$ is
\[
    (AB)_{ij}=\sum_k A_{ik}B_{kj}.
\]
Taking the differential entry by entry,
\[
    d(AB)_{ij}
    =
    \sum_k dA_{ik}\,B_{kj}+
    \sum_k A_{ik}\,dB_{kj}.
\]
This is exactly the $(i,j)$-entry of $(dA)B+A(dB)$.

For transpose, the $(i,j)$-entry of $d(A^{\top})$ is $d(A^{\top}_{ij})=d(A_{ji})$, which is the $(i,j)$-entry of $(dA)^{\top}$.
Finally,
\[
    d\tr(A)
    =d\left(\sum_i A_{ii}\right)
    =\sum_i dA_{ii}
    =\tr(dA).
\]
\end{pfbox}

Recall for a scalar-valued function $f$, the gradient is characterized by
\[
    df=\tr\bigl((\nabla_X f)^{\top} dX\bigr).
\]
This gives a concrete strategy for computing matrix gradients:
\begin{pbox}
 \begin{enumerate}
    \item Express the scalar function using traces whenever convenient.
    \item Take the differential using the rules above.
    \item Move terms cyclically inside traces until the result has the form $\tr(G^{\top} dX)$.
    \item Conclude that $\nabla_X f=G$.
\end{enumerate}   
\end{pbox}

The most frequently used trace identities are
\[
    \tr(PQ)=\tr(QP),
    \quad
    \tr(PQR)=\tr(QRP)=\tr(RPQ),
    \quad
    \tr(P)=\tr(P^{\top}).
\]
Only cyclic permutations are valid in general; the order of factors may not be
arbitrarily rearranged.

\begin{example}
Let $f(X)=\norm{X}_F^2$. Since $\norm{X}_F^2=\tr(X^{\top}X)$,
\[
\begin{aligned}
    df
    &=
    d\tr(X^{\top}X)
    =
    \tr(d(X^{\top}X))\\
    &=
    \tr((dX)^{\top}X)+\tr(X^{\top} dX).
\end{aligned}
\]
Using $\tr((dX)^{\top}X)=\tr(((dX)^{\top}X)^{\top})=\tr(X^{\top}dX)$, we get
\[
    df=2\tr(X^{\top}dX)=\tr((2X)^{\top}dX).
\]
Therefore
\[
    \nabla_X f=2X.
\]
\end{example}

\begin{example}
Let $f:M_n(\mathbb{R})\to\mathbb{R}$ be defined by $f(X)=\tr(X^2)$. Then
\[
\begin{aligned}
    df
    &=
    d\tr(X^2)
    =
    \tr(d(X\cdot X))\\
    &=
    \tr((dX)X)+\tr(X(dX))\\
    &=
    \tr(XdX)+\tr(XdX)\\
    &=
    \tr(2X\,dX).
\end{aligned}
\]
Since $df=\tr(G^{\top} dX)$, we have $G^{\top}=2X$. Therefore
\[
    \nabla_X f=2X^{\top}.
\]
\end{example}

\begin{remark}
Consider $F:M_n(\mathbb{R})\to M_n(\mathbb{R})$ defined by $F(X)=X^2$. Its differential is
\[
    dF=d(X^2)=d(X\cdot X)=(dX)X+X(dX).
\]
Thus the derivative at $X$ is the linear map
\[
    dF_X(H)=HX+XH.
\]
Since $F$ is not scalar-valued, the notion of gradient does not apply here. But its derivative and differential are still well-defined.
\end{remark}

\begin{theorem}[Chain Rule]
\label{thm:matrix-chain-rule}
Let $f:\mathcal{X}\to\mathcal{Y}$ and $g:\mathcal{Y}\to\mathcal{Z}$ be
differentiable maps between finite-dimensional matrix spaces. Then
\[
    d(g\circ f)_X
    =
    dg_{f(X)}\circ df_X.
\]
Equivalently, for every direction $H\in\mathcal{X}$,
\[
    d(g\circ f)_X(H)
    =
    dg_{f(X)}\bigl(df_X(H)\bigr).
\]
\end{theorem}

There is a particularly useful gradient form of this rule. If
$L:\mathcal{X}\to\mathcal{Y}$ is linear, its adjoint
$L^*:\mathcal{Y}\to\mathcal{X}$ is characterized by
\[
    \innerproduct{L(H)}{K}_{\mathcal{Y}}
    =
    \innerproduct{H}{L^*(K)}_{\mathcal{X}}
    \qquad
    \text{for all }H\in\mathcal{X},\ K\in\mathcal{Y}.
\]
Here $H$ is an arbitrary input direction, while $K$ is an arbitrary test
direction in the output space.
The two inner products are dictated by the spaces in which the vectors live:
$L(H)$ and $K$ belong to $\mathcal{Y}$, whereas $H$ and $L^*(K)$ belong to
$\mathcal{X}$. Thus $L^*$ depends on the chosen inner products on both
$\mathcal{X}$ and $\mathcal{Y}$; under the usual Euclidean or Frobenius inner
products, its matrix representation is the transpose of the matrix representing
$L$.

\begin{theorem}
Let $f:\mathcal{X}\to\mathcal{Y}$ and $g:\mathcal{Y}\to\mathbb{R}$ be
differentiable. Then
\[
    \nabla_X(g\circ f)
    =
    df_X^*\bigl(\nabla_Y g(f(X))\bigr),
\]
where $df_X^*$ is the adjoint of the linear map $df_X$.
\end{theorem}

\begin{pfbox}
For every input direction $H\in\mathcal{X}$, the ordinary chain rule gives
\[
    d(g\circ f)_X(H)
    =
    dg_{f(X)}\bigl(df_X(H)\bigr).
\]
Now $df_X(H)\in\mathcal{Y}$ is a direction at $f(X)$. Since
$g:\mathcal{Y}\to\mathbb{R}$ is scalar-valued, the definition of its gradient
states that, for every $Y,K\in\mathcal{Y}$,
\[
    dg_Y(K)
    =\innerproduct{\nabla_Yg(Y)}{K}_{\mathcal{Y}}
    =\innerproduct{K}{\nabla_Yg(Y)}_{\mathcal{Y}},
\]
where the second equality uses symmetry of the real Frobenius inner product.
Taking $Y=f(X)$ and $K=df_X(H)$ therefore gives
\[
    d(g\circ f)_X(H)
    =
    \innerproduct{df_X(H)}{\nabla_Yg(f(X))}_{\mathcal{Y}}.
\]
Both entries lie in $\mathcal{Y}$. To identify the gradient with respect to
$X$, however, the differential must have the form
$\innerproduct{H}{G}_{\mathcal{X}}$ for some $G\in\mathcal{X}$. The adjoint is
defined precisely to move $df_X$ across the inner product:
\[
    \innerproduct{df_X(H)}{\nabla_Y g(f(X))}_{\mathcal{Y}}
    =
    \innerproduct{H}{df_X^*(\nabla_Y g(f(X)))}_{\mathcal{X}}.
\]
Hence $G=df_X^*(\nabla_Y g(f(X)))$, which proves the result. In this sense, the
adjoint pulls the output gradient in $\mathcal{Y}$ back to an input gradient in
$\mathcal{X}$.
\end{pfbox}

\begin{remark}[Jacobian form]
Let $f:\mathbb{R}^n\to\mathbb{R}^m$ and
$g:\mathbb{R}^m\to\mathbb{R}$. Writing
$f=(f_1,\ldots,f_m)^{\top}$, the Jacobian of $f$ is
\[
    J_f(\mathbf{x})
    =
    \left(
        \frac{\partial f_i}{\partial x_j}(\mathbf{x})
    \right)_{i,j}
    \in M_{m\times n}(\mathbb{R}),
\]
and it represents the derivative:
\[
    df_{\mathbf{x}}(\mathbf{h})=J_f(\mathbf{x})\mathbf{h}.
\]
With the usual inner products, the adjoint of this derivative is represented
by the transpose Jacobian, since
\[
    \innerproduct{J_f(\mathbf{x})\mathbf{h}}{\mathbf{k}}
    =
    \innerproduct{\mathbf{h}}{J_f(\mathbf{x})^{\top}\mathbf{k}}.
\]
Therefore the gradient chain rule becomes
\[
    \nabla_{\mathbf{x}}(g\circ f)(\mathbf{x})
    =
    J_f(\mathbf{x})^{\top}
    \nabla_{\mathbf{y}}g(f(\mathbf{x})).
\]
Thus the adjoint formula is the coordinate-free version of ``multiply the
output gradient by the transpose Jacobian.''
\end{remark}

\begin{example}
Let $f:\mathbb{R}^2\to\mathbb{R}^2$ and
$g:\mathbb{R}^2\to\mathbb{R}$ be defined by
\[
    f(\mathbf{x})
    =
    \begin{pmatrix}
        x_1^2+x_2\\
        x_1x_2
    \end{pmatrix},
    \qquad
    g(\mathbf{y})=\frac{1}{2}(y_1^2+y_2^2).
\]
Then
\[
    J_f(\mathbf{x})
    =
    \begin{pmatrix}
        2x_1 & 1\\
        x_2 & x_1
    \end{pmatrix},
    \qquad
    \nabla_{\mathbf{y}}g(\mathbf{y})=\mathbf{y}.
\]
Hence
\[
\begin{aligned}
    \nabla_{\mathbf{x}}(g\circ f)(\mathbf{x})
    &=
    J_f(\mathbf{x})^{\top}f(\mathbf{x})\\
    &=
    \begin{pmatrix}
        2x_1(x_1^2+x_2)+x_1x_2^2\\
        x_1^2+x_2+x_1^2x_2
    \end{pmatrix}.
\end{aligned}
\]
This agrees with direct differentiation of
\[
    (g\circ f)(\mathbf{x})
    =
    \frac{1}{2}\left((x_1^2+x_2)^2+(x_1x_2)^2\right).
\]
\end{example}

\begin{proposition}[Affine Matrix Least Squares]
Let
\[
    A\in M_{m\times p}(\mathbb{R}),
    \qquad
    X\in M_{p\times q}(\mathbb{R}),
    \qquad
    B\in M_{q\times n}(\mathbb{R}),
    \qquad
    C\in M_{m\times n}(\mathbb{R}).
\]
Define
\[
    F(X)
    =
    \frac{1}{2}\norm{AXB-C}_F^2.
\]
Then
\[
    \nabla_X F
    =
    A^\top(AXB-C)B^\top.
\]
\label{affine least square}
\end{proposition}

\begin{pfbox}
Let $Y=AXB-C$.

Then
\[
    F(X)
    =
    \frac{1}{2}\norm{Y}_F^2
    =
    \frac{1}{2}\tr(Y^\top Y).
\]
Taking differential, we get
\[
    dY
    =
    d(AXB-C)
    =
    A(dX)B.
\]
Since $d\left(\frac{1}{2}\norm{Y}_F^2\right)=\innerproduct{Y}{dY}_F$, we have
\[
\begin{aligned}
    dF
    &=
    \innerproduct{Y}{dY}_F                                      \\[2mm]
    &=
    \innerproduct{AXB-C}{A(dX)B}_F                                \\[2mm]
    &=
    \tr\left((AXB-C)^\top A(dX)B\right).
\end{aligned}
\]
Now move the matrices cyclically inside the trace:
\[
\begin{aligned}
    dF
    &=
    \tr\left(B(AXB-C)^\top A(dX)\right)                           \\[2mm]
    &=
    \tr\left(\bigl(A^\top(AXB-C)B^\top\bigr)^\top dX\right).
\end{aligned}
\]
By the definition of the gradient,
\[
    dF
    =
    \tr\left((\nabla_X F)^\top dX\right).
\]
Hence
\[
    \nabla_X F
    =
    A^\top(AXB-C)B^\top.
\]
\end{pfbox}

\begin{remark}[Least-squares interpretations]
The rule $\nabla_X\frac12\norm{AXB-C}_F^2
=A^\top(AXB-C)B^\top$ reflects the two adjoints: left multiplication by $A$
contributes $A^\top$, while right multiplication by $B$ contributes $B^\top$.
For $B=I$, it reduces to the usual matrix least-squares gradient
$\nabla_X\frac12\norm{AX-C}_F^2=A^\top(AX-C)$.

For the linear model $y=X_{\mathrm{data}}\beta+\varepsilon$, take
$A=X_{\mathrm{data}}$, $X=\beta$, $B=I$, and $C=y$. Then
\[
    \nabla_\beta\frac12\norm{X_{\mathrm{data}}\beta-y}^2
    =X_{\mathrm{data}}^\top(X_{\mathrm{data}}\beta-y),
\]
so the first-order condition is
$X_{\mathrm{data}}^\top X_{\mathrm{data}}\hat\beta
=X_{\mathrm{data}}^\top y$. 

If
$X_{\mathrm{data}}^\top X_{\mathrm{data}}$ is invertible, equivalently if
$X_{\mathrm{data}}$ has full column rank, then
\[
    \hat\beta
    =(X_{\mathrm{data}}^\top X_{\mathrm{data}})^{-1}
      X_{\mathrm{data}}^\top y.
\]
For multiple outcomes $Y=X_{\mathrm{data}}\Gamma+E$, the same rule gives
$\nabla_\Gamma\frac12\norm{X_{\mathrm{data}}\Gamma-Y}_F^2
=X_{\mathrm{data}}^\top(X_{\mathrm{data}}\Gamma-Y)$; thus matrix least
squares is OLS applied to all outcome columns at once.

The factor $B$ allows a transformation on the right of a coefficient matrix.
For example, in
\[
    Y=X_{\mathrm{data}}\Theta R+E,
\]
$\Theta$ is the free coefficient matrix to be estimated, whereas $R$ is a
known, fixed matrix that specifies how the columns of $\Theta$ are combined.
The objective $\frac12\norm{X_{\mathrm{data}}\Theta R-Y}_F^2$ is the full
form with $A=X_{\mathrm{data}}$, $X=\Theta$, $B=R$, and $C=Y$.

For an ordinary coefficient restriction such as $\beta_1=\beta_2$, it is more
natural to parameterize the coefficient vector as $\beta=R\theta$. For example,
if $\beta\in\mathbb{R}^3$, set
\[
    R=
    \begin{bmatrix}
        1&0\\
        1&0\\
        0&1
    \end{bmatrix},
    \qquad
    \theta=
    \begin{bmatrix}
        \theta_1\\
        \theta_2
    \end{bmatrix}.
\]
Then
\[
    \beta=R\theta
    =
    \begin{bmatrix}
        \theta_1\\
        \theta_1\\
        \theta_2
    \end{bmatrix},
\]
so $\beta_1=\beta_2$ holds automatically, while $\theta_1,\theta_2$ remain
unrestricted. The restricted least-squares problem becomes
$\frac12\norm{X_{\mathrm{data}}R\theta-y}^2$, which has the same form with
$A=X_{\mathrm{data}}R$, $X=\theta$, $B=I$, and $C=y$.

Weighted least squares and GLS enter through a transformation on the left.
Suppose
\[
    y=X_{\mathrm{data}}\beta+\varepsilon,
    \qquad
    \mathbb{E}(\varepsilon)=0,
    \qquad
    \operatorname{Cov}(\varepsilon)=\Omega\succ0.
\]
GLS weights the residual by its inverse covariance and minimizes
\[
\begin{aligned}
    Q_{\mathrm{GLS}}(\beta)
    &=
    \frac12(X_{\mathrm{data}}\beta-y)^\top
    \Omega^{-1}(X_{\mathrm{data}}\beta-y)\\
    &=
    \frac12\norm{\Omega^{-1/2}(X_{\mathrm{data}}\beta-y)}^2.
\end{aligned}
\]
Thus GLS is ordinary least squares applied to the whitened data
$\widetilde X=\Omega^{-1/2}X_{\mathrm{data}}$ and
$\widetilde y=\Omega^{-1/2}y$. In the notation of $F(X)$, this corresponds to
$X=\beta$, $A=\Omega^{-1/2}X_{\mathrm{data}}$, $B=I$, and
$C=\Omega^{-1/2}y$. Therefore
\[
    \nabla_\beta Q_{\mathrm{GLS}}(\beta)
    =
    X_{\mathrm{data}}^\top\Omega^{-1}
    (X_{\mathrm{data}}\beta-y),
\]
and the first-order condition is the GLS normal equation
\[
    X_{\mathrm{data}}^\top\Omega^{-1}X_{\mathrm{data}}
    \hat\beta_{\mathrm{GLS}}
    =
    X_{\mathrm{data}}^\top\Omega^{-1}y.
\]
If $X_{\mathrm{data}}^\top\Omega^{-1}X_{\mathrm{data}}$ is invertible, then
\[
    \hat\beta_{\mathrm{GLS}}
    =
    (X_{\mathrm{data}}^\top\Omega^{-1}X_{\mathrm{data}})^{-1}
    X_{\mathrm{data}}^\top\Omega^{-1}y.
\]
Weighted least squares is the special case in which $\Omega$ is diagonal;
general GLS also permits correlated errors through the off-diagonal entries of
$\Omega$.
\end{remark}



\begin{proposition}[Derivative of Affine Mapping]
Let $A\in M_{m\times n}(\mathbb{R})$ and $\mathbf{b}\in\mathbb{R}^m$. Define
\[
    F:\mathbb{R}^n\to\mathbb{R}^m,
    \qquad
    F(\mathbf{x})=A\mathbf{x}+\mathbf{b}.
\]
Then $F$ is differentiable at every $\mathbf{x}\in\mathbb{R}^n$, and
\[
    dF_{\mathbf{x}}(\mathbf{h})=A\mathbf{h}.
\]
Thus the matrix representing the derivative $dF_{\mathbf{x}}$ is $A$.
\end{proposition}

\begin{pfbox}
Consider the candidate linear map $L:\mathbb{R}^n\to\mathbb{R}^m$ defined by
$L(\mathbf{h})=A\mathbf{h}$. For every nonzero $\mathbf{h}$,
\[
\begin{aligned}
    &\frac{\norm{F(\mathbf{x}+\mathbf{h})-F(\mathbf{x})-L(\mathbf{h})}}
    {\norm{\mathbf{h}}}\\
    =& \frac{\norm{A(\mathbf{x}+\mathbf{h})+\mathbf{b}
    -(A\mathbf{x}+\mathbf{b})-A\mathbf{h}}}{\norm{\mathbf{h}}}
    =0.
\end{aligned}
\]
The quotient is identically zero, so its limit as $\mathbf{h}\to\mathbf{0}$ is
zero. By the definition of the Fr\'echet derivative,
$dF_{\mathbf{x}}=L$, whose matrix is $A$.
\end{pfbox}

\begin{remark}
Because $F$ is vector-valued, $A$ is its derivative (or Jacobian), not its
gradient. Componentwise, the rows of $A$ are
$(\nabla_{\mathbf{x}}F_i)^{\top}$.
\end{remark}


\begin{theorem}
Let $A\in M_n(\mathbb{R})$ and define $f:\mathbb{R}^n\to\mathbb{R}$ by
$f(\mathbf{x})=\mathbf{x}^{\top} A\mathbf{x}$. Then
\[
    \nabla_{\mathbf{x}} f=(A+A^{\top})\mathbf{x}.
\]
In particular, if $A$ is symmetric, then $\nabla_{\mathbf{x}} f=2A\mathbf{x}$.
\end{theorem}

\begin{pfbox}
The matrix $A$ is not assumed to be symmetric. Since $\mathbf{x}^{\top} A\mathbf{x}$ is a scalar,
\[
    f(\mathbf{x})=\tr(\mathbf{x}^{\top} A\mathbf{x}).
\]
Taking the differential,
\[
\begin{aligned}
    df
    &=
    \tr(d(\mathbf{x}^{\top} A\mathbf{x}))\\
    &=
    \tr((d\mathbf{x})^{\top} A\mathbf{x})+\tr(\mathbf{x}^{\top} A\,d\mathbf{x}).
\end{aligned}
\]
For the first term,
\[
    \tr((d\mathbf{x})^{\top} A\mathbf{x})
    =
    \tr(((d\mathbf{x})^{\top} A\mathbf{x})^{\top})
    =
    \tr(\mathbf{x}^{\top} A^{\top} d\mathbf{x}).
\]
Hence
\[
\begin{aligned}
    df
    &=
    \tr(\mathbf{x}^{\top} A^{\top} d\mathbf{x})+\tr(\mathbf{x}^{\top} A\,d\mathbf{x})\\
    &=
    \tr\bigl(\mathbf{x}^{\top}(A^{\top}+A)d\mathbf{x}\bigr).
\end{aligned}
\]
Comparing this with $df=\tr((\nabla_{\mathbf{x}} f)^{\top} d\mathbf{x})$ gives
\[
    (\nabla_{\mathbf{x}} f)^{\top}=\mathbf{x}^{\top}(A^{\top}+A).
\]
Therefore
\[
    \nabla_{\mathbf{x}} f=(A+A^{\top})\mathbf{x}.
\]
If $A=A^{\top}$, this becomes $\nabla_{\mathbf{x}} f=2A\mathbf{x}$.
\end{pfbox}



\begin{theorem}
In each item below, all matrices and vectors other than the displayed variables
are constant.
\begin{enumerate}
    \item If $\mathbf{x}\in\mathbb{R}^n$,
    $A\in M_{m\times n}(\mathbb{R})$, and
    $B\in M_{n\times m}(\mathbb{R})$, then
    \[
        \frac{\partial(A\mathbf{x})}{\partial\mathbf{x}^{\top}}=A,
        \qquad
        \frac{\partial(\mathbf{x}^{\top}B)}{\partial\mathbf{x}}=B.
    \]

    \item If $\mathbf{x},\mathbf{a}\in\mathbb{R}^n$, then
    $\displaystyle
    \frac{\partial}{\partial\mathbf{x}}(\mathbf{a}^{\top}\mathbf{x})
    =\frac{\partial}{\partial\mathbf{x}}(\mathbf{x}^{\top}\mathbf{a})
    =\mathbf{a}$.

    \item If $\mathbf{x}\in\mathbb{R}^m$,
    $\mathbf{y}\in\mathbb{R}^n$, and
    $A\in M_{m\times n}(\mathbb{R})$, define
    $f(\mathbf{x},\mathbf{y})=\mathbf{x}^{\top}A\mathbf{y}\in\mathbb{R}$.
    Then
    $\frac{\partial f}{\partial\mathbf{x}}=A\mathbf{y}$
    and $\frac{\partial f}{\partial\mathbf{y}}=A^{\top}\mathbf{x}$.

    \item If $\mathbf{x}\in\mathbb{R}^n$ and $A\in M_n(\mathbb{R})$, then
    $\displaystyle
    \frac{\partial}{\partial\mathbf{x}}(\mathbf{x}^{\top}A\mathbf{x})
    =(A+A^{\top})\mathbf{x}$ and
    $\displaystyle
    \frac{\partial^2}{\partial\mathbf{x}\partial\mathbf{x}^{\top}}
    (\mathbf{x}^{\top}A\mathbf{x})=A+A^{\top}$.

    \item If $X,A\in M_{m\times n}(\mathbb{R})$, then
    $\displaystyle \frac{\partial}{\partial X}\tr(A^{\top}X)=A$.

    \item If $X\in M_{m\times n}(\mathbb{R})$, then
    $\displaystyle
    \frac{\partial}{\partial X}\left(\frac{1}{2}\norm{X}_F^2\right)=X$.

    \item If $X\in M_{m\times n}(\mathbb{R})$,
    $A\in M_{p\times m}(\mathbb{R})$,
    $B\in M_{n\times q}(\mathbb{R})$, and
    $C\in M_{p\times q}(\mathbb{R})$, then
    $\displaystyle
    \frac{\partial}{\partial X}\left(\frac{1}{2}\norm{AXB-C}_F^2\right)
    =A^{\top}(AXB-C)B^{\top}$.

    % \item If $X\in M_{m\times n}(\mathbb{R})$, $A\in M_m(\mathbb{R})$,
    % and $B\in M_n(\mathbb{R})$, then
    % \[
    %     \frac{\partial}{\partial X}\tr(X^{\top}AXB)
    %     =AXB+A^{\top}XB^{\top}.
    % \]
\end{enumerate}
\end{theorem}

\begin{pfbox}
For (1), define
\[
    F(\mathbf{x})=A\mathbf{x},
    \qquad
    G(\mathbf{x})=\mathbf{x}^{\top}B.
\]
The preceding proposition gives
$\partial F/\partial\mathbf{x}^{\top}=A$. Moreover,
$G^{\top}(\mathbf{x})=B^{\top}\mathbf{x}$, so
\[
    \frac{\partial G}{\partial\mathbf{x}}
    =
    \left(
        \frac{\partial G^{\top}}{\partial\mathbf{x}^{\top}}
    \right)^{\top}
    =(B^{\top})^{\top}
    =B.
\]
For (2), let
$f(\mathbf{x})=\mathbf{a}^{\top}\mathbf{x}=\mathbf{x}^{\top}\mathbf{a}$. Then
\[
    df=\tr(\mathbf{a}^{\top}d\mathbf{x}),
\]
so $\frac{\partial f}{\partial\mathbf{x}}
=(\mathbf{a}^{\top})^{\top}=\mathbf{a}$.

For (3), the product rule gives
\[
    df
    =
    (d\mathbf{x})^{\top}A\mathbf{y}
    +\mathbf{x}^{\top}A\,d\mathbf{y}
    =
    (A\mathbf{y})^{\top}d\mathbf{x}
    +(A^{\top}\mathbf{x})^{\top}d\mathbf{y}.
\]
Since $\mathbf{x}$ and $\mathbf{y}$ are both column variables, their total
differential has the standard form
\[
    df
    =
    \left(\frac{\partial f}{\partial\mathbf{x}}\right)^{\top}d\mathbf{x}
    +
    \left(\frac{\partial f}{\partial\mathbf{y}}\right)^{\top}d\mathbf{y}.
\]
Comparing the coefficients of $d\mathbf{x}$ and $d\mathbf{y}$ identifies both
vector partial derivatives. 

For (5), let $f(X)=\tr(A^{\top}X)$. Then
\[
    df=\tr(A^{\top}dX),
\]
so $\frac{\partial f}{\partial X}=A$. The equivalent formula follows by
replacing $A^{\top}$ with $A$.

% For (8), use the product rule and cyclicity of trace:
% \[
% \begin{aligned}
%     d\tr(X^{\top}AXB)
%     &=\tr((dX)^{\top}AXB)+\tr(X^{\top}A\,dX B)\\
%     &=\tr\left((AXB+A^{\top}XB^{\top})^{\top}dX\right).
% \end{aligned}
% \]
% Therefore
% \[
%     \frac{\partial}{\partial X}\tr(X^{\top}AXB)
%     =AXB+A^{\top}XB^{\top}.
% \]
Parts (4), (6), (7) were proved earlier.
Note that every function in parts (2)--(7) is scalar-valued. Under our column convention,
the vector or matrix obtained by collecting its partial derivatives is exactly
its gradient; that is,
\[
    \frac{\partial f}{\partial\mathbf{x}}=\nabla_{\mathbf{x}}f,
    \qquad
    \frac{\partial f}{\partial\mathbf{y}}=\nabla_{\mathbf{y}}f,
    \qquad
    \frac{\partial f}{\partial X}=\nabla_Xf.
\]
\end{pfbox}

% \begin{remark}
% A simple application of (8) is multivariate linear regression. For observation $i$,
% let $\mathbf{z}_i\in\mathbb{R}^k$ contain its $k$ regressors and let
% $\mathbf{y}_i\in\mathbb{R}^q$ contain its $q$ outcomes. The model is
% \[
%     \mathbf{y}_i
%     =
%     \Gamma^{\top}\mathbf{z}_i+\mathbf{e}_i,
%     \qquad
%     \mathbb{E}(\mathbf{e}_i)=0,
%     \qquad
%     \operatorname{Cov}(\mathbf{e}_i)=\Sigma\succ0,
% \]
% where $\Gamma\in M_{k\times q}(\mathbb{R})$ is the coefficient matrix and
% $\Sigma\in M_q(\mathbb{R})$ is the covariance matrix of the $q$ errors. Stacking the $n$ observations gives $Y=Z\Gamma+E$.

% Here $Y\in M_{n\times q}(\mathbb{R})$ has row $\mathbf{y}_i^{\top}$,
% $E\in M_{n\times q}(\mathbb{R})$ has row $\mathbf{e}_i^{\top}$, and the
% \emph{design matrix} $Z\in M_{n\times k}(\mathbb{R})$ has row
% $\mathbf{z}_i^{\top}$. Thus each row of $Z$ records the regressors for one
% observation; its first column is usually all ones when the model has an
% intercept. Each column of $\Gamma$ contains the coefficients for one outcome.

% To estimate $\Gamma$, we minimize the residual criterion, i.e. the loss function
% measuring how far the fitted value $Z\Gamma$ is from the observed outcome $Y$:
% \[
%     Q(\Gamma)
%     =
%     \frac12\tr\left(
%         (Y-Z\Gamma)^{\top}(Y-Z\Gamma)\Sigma^{-1}
%     \right).
% \]
% Writing it out,
% \[
% \begin{aligned}
%     Q(\Gamma)
%     &=
%     \frac12\tr(Y^{\top}Y\Sigma^{-1})
%     -\frac12\tr(Y^{\top}Z\Gamma\Sigma^{-1})
%     -\frac12\tr(\Gamma^{\top}Z^{\top}Y\Sigma^{-1}) \\
%     &\qquad
%     +\underbrace{
%         \frac12\tr\left(\Gamma^{\top}(Z^{\top}Z)\Gamma\Sigma^{-1}\right)
%     }_{\text{the quadratic term in }\Gamma}.
% \end{aligned}
% \]
% Only the last term uses item (8):
% \[
% \begin{aligned}
%     \nabla_{\Gamma}
%     \frac12\tr\left(\Gamma^{\top}(Z^{\top}Z)\Gamma\Sigma^{-1}\right)
%     &=
%     \frac12\left[
%         (Z^{\top}Z)\Gamma\Sigma^{-1}
%         +(Z^{\top}Z)^{\top}\Gamma(\Sigma^{-1})^{\top}
%     \right]  \\
%     &=
%     Z^{\top}Z\Gamma\Sigma^{-1},
% \end{aligned}
% \]
% because $Z^{\top}Z$ and $\Sigma^{-1}$ are symmetric. The first term is constant
% in $\Gamma$, and the two linear terms give $-Z^{\top}Y\Sigma^{-1}$. Hence
% \[
%     \nabla_{\Gamma}Q
%     =
%     Z^{\top}(Z\Gamma-Y)\Sigma^{-1}.
% \]
% Since $\Sigma^{-1}$ is invertible, setting the gradient equal to zero gives
% $Z^{\top}Z\widehat\Gamma=Z^{\top}Y$. If $Z$ has full column rank, then
% \[
%     \widehat\Gamma
%     =
%     (Z^{\top}Z)^{-1}Z^{\top}Y.
% \]
% Thus this is ordinary OLS applied separately to each outcome column. The
% covariance $\Sigma$ affects the uncertainty of the estimates, but it cancels
% from the coefficient formula because every outcome uses the same design matrix
% $Z$.
% \end{remark}


\begin{remark}[Derivative-layout convention]
The transpose records the \emph{layout} of the derivative. For a scalar-valued function
$f:\mathbb{R}^n\to\mathbb{R}$, we use the column-gradient convention
\[
    df
    =
    \left(\frac{\partial f}{\partial\mathbf{x}}\right)^{\top}d\mathbf{x},
    \qquad
    \frac{\partial f}{\partial\mathbf{x}}\in\mathbb{R}^n.
\]
For a column-vector-valued function
$F:\mathbb{R}^n\to\mathbb{R}^m$, its derivative must instead be an
$m\times n$ matrix that acts on $d\mathbf{x}$:
\[
    dF
    =
    \underbrace{\frac{\partial F}{\partial\mathbf{x}^{\top}}}_{J_F\in M_{m\times n}(\mathbb{R})}
    d\mathbf{x}.
\]
Thus, $\partial F/\partial\mathbf{x}^{\top}$ is the usual $m\times n$
Jacobian, whose rows correspond to the components $F_i$ and whose columns
correspond to the variables $x_j$. For a row-vector-valued function
$G:\mathbb{R}^n\to\mathbb{R}^{1\times m}$, we instead define
\[
    dG
    =
    d\mathbf{x}^{\top}
    \underbrace{\frac{\partial G}{\partial\mathbf{x}}}_{\in M_{n\times m}(\mathbb{R})},
\]
so that the output remains a row vector. Consequently,
\[
    d(A\mathbf{x})=A\,d\mathbf{x}
    \quad\Longrightarrow\quad
    \frac{\partial(A\mathbf{x})}{\partial\mathbf{x}^{\top}}=A,
    \qquad
    d(\mathbf{x}^{\top}B)=d\mathbf{x}^{\top}B
    \quad\Longrightarrow\quad
    \frac{\partial(\mathbf{x}^{\top}B)}{\partial\mathbf{x}}=B.
\]
This is the layout convention used throughout this section. Equivalently, one
may first transpose the row-valued function and compute the ordinary Jacobian
of $G^{\top}=B^{\top}\mathbf{x}$, then transpose the result back to the
row-output layout:
\[
    \frac{\partial G}{\partial\mathbf{x}}
    =
    \left(
        \frac{\partial G^{\top}}{\partial\mathbf{x}^{\top}}
    \right)^{\top}
    =
    (B^{\top})^{\top}
    =B.
\]
\end{remark}

\subsubsection*{Invertible matrices}

The following formulas apply to functions whose domain is the set of invertible matrices
\[
    GL_n(\mathbb{R})=\{X\in M_n(\mathbb{R}):\det X\ne 0\},
\]
called the general linear group over $\mathbb{R}$.

\begin{theorem}[Inversion]
The inversion map $X\mapsto X^{-1}$ on $GL_n(\mathbb{R})$ satisfies
\[
    d(X^{-1})=-X^{-1}(dX)X^{-1}.
\]
\end{theorem}

\begin{pfbox}
Start from the identity $X^{-1}X=I$. Differentiating both sides gives
$d(X^{-1})\,X+X^{-1}dX=0$, hence $d(X^{-1})\,X=-X^{-1}dX$.
Multiplying on the right by $X^{-1}$ gives
\[
    d(X^{-1})=-X^{-1}(dX)X^{-1}.
\]
The inversion map is matrix-valued, not scalar-valued, so we talk about its derivative or differential, not its gradient.
\end{pfbox}

Set $GL_n^+(\mathbb{R})=\{X\in GL_n(\mathbb{R}):\det X>0\}$.

\begin{theorem}[Log-determinant]
The function $f:GL_n^+(\mathbb{R})\to\mathbb{R}$ defined by $f(X)=\log\det X$ has gradient
\[
    \nabla_X(\log\det X)=(X^{-1})^{\top}.
\]
Equivalently,
\[
    d(\log\det X)=\tr(X^{-1}dX).
\]
\end{theorem}

\begin{pfbox}
Let $f(X)=\log\det X$. For a small matrix $H$,
\[
\begin{aligned}
    f(X+H)-f(X)
    &=
    \log\det(X+H)-\log\det X\\
    &=
    \log\det\bigl(X(I+X^{-1}H)\bigr)-\log\det X\\
    &=
    \log\det X+
    \log\det(I+X^{-1}H)-\log\det X\\
    &=
    \log\det(I+X^{-1}H).
\end{aligned}
\]
Now put $K=X^{-1}H$. We claim that
\[
    \det(I+K)=1+\tr(K)+O(\norm{K}_F^2).
\]
Let $\operatorname{Spec}(K)=\{\lambda_1,\ldots,\lambda_n\}$, counted
with algebraic multiplicity. Then
\[
\begin{aligned}
    \det(I+K)
    &=
    \prod_{i=1}^n(1+\lambda_i)                                  \\
    &=
    1+\sum_{i=1}^n\lambda_i+\text{higher-order terms}             \\
    &=
    1+\tr(K)+O(\norm{K}_F^2).
\end{aligned}
\]
The higher-order terms are sums of products of at least two eigenvalues, hence
are $O(\norm{K}_F^2)$ since $|\lambda_i|\le\norm{K}_F\le\norm{X^{-1}}\norm{H}_F\to 0$ as $H\to 0$.

Next write $u=\tr(K)+O(\norm{K}_F^2)$. The one-variable Taylor expansion of
$g(u)=\log(1+u)$ at $u=0$ gives $g(0)=0$ and $g'(0)=1$, so
$\log(1+u)=u+O(u^2)$ as $u\to0$. Hence
\[
\begin{aligned}
    \log\det(I+K)
    &=
    \log\bigl(1+\tr(K)+O(\norm{K}_F^2)\bigr)\\
    &= 
    \tr(K)+O(\norm{K}_F^2),
\end{aligned}
\]
where $O(\norm{K}_F^2)=O(\norm{H}_F^2)=o(\norm{H}_F)$. Therefore
\[
    f(X+H)-f(X)
    =
    \tr(X^{-1}H)+o(\norm{H}_F).
\]
In differential notation, $df=\tr(X^{-1}dX)$, so $\nabla_X f=(X^{-1})^{\top}$.
\end{pfbox}

\begin{corollary}[Determinant]
On $GL_n^+(\mathbb{R})$,
$d(\det X)=\det(X)\tr(X^{-1}dX),\quad\nabla_X\det X=\det(X)X^{-\top}$.
\end{corollary}

\begin{pfbox}
    Chain Rule.
\end{pfbox}

\subsubsection*{Hessian}

The Hessian describes how the gradient changes. In one variable, $f'(x)$
measures slope, while $f''(x)$ measures how the slope changes. In several
variables, the same idea remains true: the gradient gives the first-order
change of the function, and the Hessian gives the first-order change of the
gradient.

\paragraph*{Vector Variable}

Let $U\subseteq\mathbb{R}^n$ be open and let $f:U\to\mathbb{R}$ be
differentiable. For $\mathbf{x}\in U$, the derivative of $f$ at $\mathbf{x}$ is
the linear map $df_{\mathbf{x}}:\mathbb{R}^n\to\mathbb{R}$ characterized by
\[
    f(\mathbf{x}+\mathbf{h})-f(\mathbf{x})
    =
    df_{\mathbf{x}}(\mathbf{h})
    +
    o(\norm{\mathbf{h}}).
\]
Since $df_{\mathbf{x}}$ is a linear functional, there exists a unique vector
$\nabla f(\mathbf{x})\in\mathbb{R}^n$ such that
\[
    df_{\mathbf{x}}(\mathbf{h})
    =
    \innerproduct{\nabla f(\mathbf{x})}{\mathbf{h}}.
\]
Thus the first-order approximation of $f$ is
\[
    f(\mathbf{x}+\mathbf{h})
    =
    f(\mathbf{x})
    +
    \innerproduct{\nabla f(\mathbf{x})}{\mathbf{h}}
    +
    o(\norm{\mathbf{h}}).
\]



Now assume $f$ is $C^2$. Then the gradient $\nabla f:U\to\mathbb{R}^n$ is itself
a differentiable vector-valued function. Its derivative at $\mathbf{x}$ is the
Hessian matrix
\[
    H_f(\mathbf{x})
    =
    \left(
        \frac{\partial^2 f}{\partial x_i\partial x_j}(\mathbf{x})
    \right)_{i,j}.
\]
This means that if we perturb $\mathbf{x}$ in the direction $\mathbf{h}$ by
replacing it with $\mathbf{x}+t\mathbf{h}$, then
\[
    \nabla f(\mathbf{x}+t\mathbf{h})
    =
    \nabla f(\mathbf{x})
    +
    tH_f(\mathbf{x})\mathbf{h}
    +
    o(t).
\]

\begin{pfbox}
Let $G(\mathbf{x})=\nabla f(\mathbf{x})$. Since $G$ is differentiable at
$\mathbf{x}$, we have
\[
    G(\mathbf{x}+\mathbf{k})
    =
    G(\mathbf{x})
    +
    dG_{\mathbf{x}}(\mathbf{k})
    +
    o(\norm{\mathbf{k}}).
\]
For $G=\nabla f$, the derivative $dG_{\mathbf{x}}$ is represented by the Hessian
matrix. Hence $dG_{\mathbf{x}}(\mathbf{k})=H_f(\mathbf{x})\mathbf{k}$, so
\[
    \nabla f(\mathbf{x}+\mathbf{k})
    =
    \nabla f(\mathbf{x})
    +
    H_f(\mathbf{x})\mathbf{k}
    +
    o(\norm{\mathbf{k}}).
\]
Now set $\mathbf{k}=t\mathbf{h}$. Since
$\norm{\mathbf{k}}=\norm{t\mathbf{h}}=|t|\norm{\mathbf{h}}$, for fixed
$\mathbf{h}$ the remainder $o(\norm{t\mathbf{h}})$ can be written as $o(t)$.
Therefore
\[
    \nabla f(\mathbf{x}+t\mathbf{h})
    =
    \nabla f(\mathbf{x})
    +
    tH_f(\mathbf{x})\mathbf{h}
    +
    o(t).
\]
\end{pfbox}

\begin{definition}[Vector Hessian]
Let $f:U\to\mathbb{R}$ be $C^2$. The Hessian of $f$ at $\mathbf{x}$ is the linear
map $\nabla^2 f(\mathbf{x}):\mathbb{R}^n\to\mathbb{R}^n$ defined by
\[
    \nabla^2 f(\mathbf{x})[\mathbf{h}]
    =
    H_f(\mathbf{x})\mathbf{h}.
\]
In coordinates, this is the usual Hessian matrix
\[
    H_f(\mathbf{x})
    =
    \left(
        \frac{\partial^2 f}{\partial x_i\partial x_j}(\mathbf{x})
    \right)_{i,j}.
\]
\end{definition}

To understand the Hessian as curvature, fix a point $\mathbf{x}$ and a direction
$\mathbf{h}$. Define the one-variable function $\phi(t)=f(\mathbf{x}+t\mathbf{h})$.
Then $\phi$ records the behavior of $f$ along the line through $\mathbf{x}$ in
the direction $\mathbf{h}$.

By the chain rule, $\phi'(t)=\innerproduct{\nabla f(\mathbf{x}+t\mathbf{h})}{\mathbf{h}}$.
Using the first-order approximation of the gradient, write
\[
    \nabla f(\mathbf{x}+t\mathbf{h})
    =
    \nabla f(\mathbf{x})
    +
    tH_f(\mathbf{x})\mathbf{h}
    +
    \mathbf{r}(t),
    \qquad
    \frac{\norm{\mathbf{r}(t)}}{|t|}\to0.
\]
Then
\[
\begin{aligned}
    \phi'(t)
    &=
    \innerproduct{
        \nabla f(\mathbf{x})
        +
        tH_f(\mathbf{x})\mathbf{h}
        +
        \mathbf{r}(t)
    }{\mathbf{h}}                                                \\[1mm]
    &=
    \innerproduct{\nabla f(\mathbf{x})}{\mathbf{h}}
    +
    t\innerproduct{H_f(\mathbf{x})\mathbf{h}}{\mathbf{h}}
    +
    \innerproduct{\mathbf{r}(t)}{\mathbf{h}}                     \\[1mm]
    &=
    \innerproduct{\nabla f(\mathbf{x})}{\mathbf{h}}
    +
    t\innerproduct{H_f(\mathbf{x})\mathbf{h}}{\mathbf{h}}
    +
    o(t),
\end{aligned}
\]
Indeed, by Cauchy--Schwarz,
\[
    \frac{
        \left|\innerproduct{\mathbf{r}(t)}{\mathbf{h}}\right|
    }{|t|}
    \le
    \frac{\norm{\mathbf{r}(t)}}{|t|}\norm{\mathbf{h}}.
\]
Since $\mathbf{h}$ is fixed and $\norm{\mathbf{r}(t)}/|t|\to0$, we get
$\innerproduct{\mathbf{r}(t)}{\mathbf{h}}=o(t)$.
Therefore
\[
    \phi''(0)
    =
    \innerproduct{H_f(\mathbf{x})\mathbf{h}}{\mathbf{h}}
    =
    \mathbf{h}^{\top}H_f(\mathbf{x})\mathbf{h}.
\]
This number measures the curvature of $f$ at $\mathbf{x}$ along the direction
$\mathbf{h}$.



\begin{remark}
The scalar
\[
    q_{\mathbf{x}}(\mathbf{h})
    =
    \mathbf{h}^{\top}H_f(\mathbf{x})\mathbf{h}
\]
is the quadratic form associated with the Hessian at $\mathbf{x}$. Its sign is
the second-order curvature of $f$ along the direction $\mathbf{h}$: positive
means upward, negative means downward, and zero means flat to second order.
\end{remark}

\begin{proposition}[Second-Order Approximation]
Let $f:U\to\mathbb{R}$ be $C^2$. Then for small $\mathbf{h}$,
\[
    f(\mathbf{x}+\mathbf{h})
    =
    f(\mathbf{x})
    +
    \innerproduct{\nabla f(\mathbf{x})}{\mathbf{h}}
    +
    \frac12\mathbf{h}^{\top}H_f(\mathbf{x})\mathbf{h}
    +
    o(\norm{\mathbf{h}}^2).
\]
\end{proposition}

\begin{pfbox}
Define $\phi(t)=f(\mathbf{x}+t\mathbf{h})$. By the one-variable Taylor expansion
at $t=0$,
\[
    \phi(1)
    =
    \phi(0)+\phi'(0)+\frac12\phi''(0)+o(\norm{\mathbf{h}}^2).
\]
Now $\phi(0)=f(\mathbf{x})$ and $\phi(1)=f(\mathbf{x}+\mathbf{h})$. Also,
\[
    \phi'(0)
    =
    \innerproduct{\nabla f(\mathbf{x})}{\mathbf{h}},
    \qquad
    \phi''(0)
    =
    \mathbf{h}^{\top}H_f(\mathbf{x})\mathbf{h}.
\]
Substituting these identities gives
\[
    f(\mathbf{x}+\mathbf{h})
    =
    f(\mathbf{x})
    +
    \innerproduct{\nabla f(\mathbf{x})}{\mathbf{h}}
    +
    \frac12\mathbf{h}^{\top}H_f(\mathbf{x})\mathbf{h}
    +
    o(\norm{\mathbf{h}}^2).
\]
\end{pfbox}


\paragraph*{Matrix Variable}

Now let $f:M_{m\times n}(\mathbb{R})\to\mathbb{R}$ be a scalar-valued function
of a matrix variable $X$. 

For a small perturbation $H\in M_{m\times n}(\mathbb{R})$, differentiability
means
\[
    f(X+H)-f(X)
    =
    df_X(H)+o(\norm{H}_F).
\]
Since $df_X$ is a linear functional on the matrix space, there exists a unique
matrix $\nabla_X f(X)$ such that
\[
    df_X(H)
    =
    \innerproduct{\nabla_X f(X)}{H}_F
    =
    \operatorname{tr}\left((\nabla_Xf(X))^\top H\right).
\]
Therefore
\[
    f(X+H)
    =
    f(X)
    +
    \innerproduct{\nabla_X f(X)}{H}_F
    +
    o(\norm{H}_F).
\]

Now assume $f$ is twice differentiable. The gradient
$\nabla_X f:M_{m\times n}(\mathbb{R})\to M_{m\times n}(\mathbb{R})$ is a
matrix-valued function. Therefore its derivative is a linear map from matrix
directions to matrix directions.

\begin{definition}[Matrix Hessian]
The Hessian of $f$ at $X$ is the linear operator
\[
    \nabla_X^2 f(X):
    M_{m\times n}(\mathbb{R})\to M_{m\times n}(\mathbb{R})
\]
defined by
\[
    \nabla_X^2 f(X)[H]
    =
    d(\nabla_X f)_X(H).
\]
In words, $\nabla_X^2f(X)[H]$ is the first-order change of the gradient when
$X$ moves in the direction $H$.
\end{definition}

Thus, if we replace $X$ by $X+tH$, then
\[
    \nabla_X f(X+tH)
    =
    \nabla_X f(X)
    +
    t\nabla_X^2 f(X)[H]
    +
    o(t).
\]

\begin{pfbox}
Let $G(X)=\nabla_X f(X)$. Since $G$ is differentiable at $X$, we have
\[
    G(X+K)
    =
    G(X)
    +
    dG_X(K)
    +
    R(K),
\]
where $\norm{R(K)}_F/\norm{K}_F\to 0$ as $K\to 0$.
By definition, $dG_X(K)=\nabla_X^2 f(X)[K]$.

Now take $K=tH$. Since $\norm{K}_F=\norm{tH}_F=|t|\norm{H}_F$, for fixed $H$ the
remainder satisfies $\norm{R(tH)}_F/|t|\to 0$. Hence $R(tH)=o(t)$, and therefore
\[
    \nabla_X f(X+tH)
    =
    \nabla_X f(X)
    +
    t\nabla_X^2 f(X)[H]
    +
    o(t).
\]
\end{pfbox}

To interpret the matrix Hessian as curvature, fix a matrix direction $H$ and
define the one-variable function $\phi(t)=f(X+tH)$. Then
$\phi'(t)=\innerproduct{\nabla_X f(X+tH)}{H}_F$. Using the first-order
approximation of the gradient,
\[
\begin{aligned}
    \phi'(t)
    &=
    \innerproduct{
        \nabla_X f(X)
        +
        t\nabla_X^2f(X)[H]
        +
        o(t)
    }{H}_F                                                       \\[1mm]
    &=
    \innerproduct{\nabla_X f(X)}{H}_F
    +
    t\innerproduct{\nabla_X^2f(X)[H]}{H}_F
    +
    o(t).
\end{aligned}
\]
Therefore
\[
    \phi''(0)
    =
    \innerproduct{\nabla_X^2f(X)[H]}{H}_F.
\]
This is the matrix-variable analogue of
$\mathbf{h}^\top H_f(\mathbf{x})\mathbf{h}$.



\begin{proposition}[Second-Order Approximation for Matrix Variables]
Let $f:M_{m\times n}(\mathbb{R})\to\mathbb{R}$ be twice differentiable. Then
\[
    f(X+H)
    =
    f(X)
    +
    \innerproduct{\nabla_X f(X)}{H}_F
    +
    \frac12\innerproduct{\nabla_X^2f(X)[H]}{H}_F
    +
    o(\norm{H}_F^2).
\]
\end{proposition}

\begin{pfbox}
Define $\phi(t)=f(X+tH)$. By the one-variable Taylor expansion,
\[
    \phi(1)
    =
    \phi(0)+\phi'(0)+\frac12\phi''(0)+o(\norm{H}_F^2).
\]
Now $\phi(0)=f(X)$ and $\phi(1)=f(X+H)$. Moreover,
\[
    \phi'(0)
    =
    \innerproduct{\nabla_X f(X)}{H}_F,
    \qquad
    \phi''(0)
    =
    \innerproduct{\nabla_X^2f(X)[H]}{H}_F.
\]
Substituting these into the Taylor expansion gives
\[
    f(X+H)
    =
    f(X)
    +
    \innerproduct{\nabla_X f(X)}{H}_F
    +
    \frac12\innerproduct{\nabla_X^2f(X)[H]}{H}_F
    +
    o(\norm{H}_F^2).
\]
\end{pfbox}

\begin{proposition}[Hessian of Affine Matrix Least Squares]
Let
\[
    f(X)
    =
    \frac12\norm{AXB-C}_F^2.
\]
Then
\[
    \nabla_X f(X)
    =
    A^\top(AXB-C)B^\top,
\]
and
\[
    \nabla_X^2 f(X)[H]
    =
    A^\top A H B B^\top.
\]
Moreover,
\[
    \innerproduct{\nabla_X^2 f(X)[H]}{H}_F
    =
    \norm{AHB}_F^2
    \ge 0.
\]
Therefore $f$ is convex in $X$.
\end{proposition}

\begin{pfbox}
Let $R(X)=AXB-C$. Then $f(X)=\frac12\norm{R(X)}_F^2$. We already know that
\[
    \nabla_X f(X)
    =
    A^\top(AXB-C)B^\top.
\]
To find the Hessian, we differentiate the gradient.

Replace $X$ by $X+tH$. Then
\[
\begin{aligned}
    \nabla_X f(X+tH)
    &=
    A^\top\bigl(A(X+tH)B-C\bigr)B^\top                         \\[1mm]
    &=
    A^\top(AXB-C)B^\top
    +
    tA^\top AHB B^\top                                          \\[1mm]
    &=
    \nabla_X f(X)
    +
    tA^\top AHB B^\top.
\end{aligned}
\]
Comparing this with
\[
    \nabla_X f(X+tH)
    =
    \nabla_X f(X)
    +
    t\nabla_X^2 f(X)[H]
    +
    o(t),
\]
we obtain
\[
    \nabla_X^2 f(X)[H]
    =
    A^\top AHB B^\top.
\]

Now compute the curvature in the direction $H$:
\[
\begin{aligned}
    \innerproduct{\nabla_X^2 f(X)[H]}{H}_F
    &=
    \innerproduct{A^\top AHB B^\top}{H}_F                       \\[1mm]
    &=
    \operatorname{tr}\left((A^\top AHB B^\top)^\top H\right)     \\[1mm]
    &=
    \operatorname{tr}\left(BB^\top H^\top A^\top AH\right)       \\[1mm]
    &=
    \operatorname{tr}\left(B^\top H^\top A^\top AHB\right)        \\[1mm]
    &=
    \operatorname{tr}\left((AHB)^\top(AHB)\right)                \\[1mm]
    &=
    \norm{AHB}_F^2
    \ge 0.
\end{aligned}
\]
Hence the curvature is nonnegative in every direction $H$. Therefore the
least-squares objective is convex in $X$.
\end{pfbox}

\begin{remark}
This proposition is important because it turns a least-squares problem into a
curvature statement. The Hessian is the linear operator
\[
    L(H)=A^\top AHB B^\top.
\]
Its quadratic form is
\[
    \innerproduct{L(H)}{H}_F=\norm{AHB}_F^2\ge0,
\]
so the least-squares objective is convex. Hence solving
$\nabla_X f(X)=0$ gives a global minimizer, not merely a local candidate. If
$AHB=\mathbf O$ only when $H=\mathbf O$, then the curvature is positive in
every nonzero direction, so the minimizer is unique.
\end{remark}

\begin{gbox}
The main points of this section are:
\begin{enumerate}
    \item The Hessian is the derivative of the gradient. For vectors,
    \[
        \nabla f(\mathbf{x}+t\mathbf{h})
        =
        \nabla f(\mathbf{x})
        +
        tH_f(\mathbf{x})\mathbf{h}
        +
        o(t).
    \]
    For matrices,
    \[
        \nabla_X f(X+tH)
        =
        \nabla_X f(X)
        +
        t\nabla_X^2f(X)[H]
        +
        o(t).
    \]

    \item Curvature is obtained by pairing the Hessian output with the original
    direction:
    \[
        \mathbf{h}^\top H_f(\mathbf{x})\mathbf{h},
        \qquad
        \innerproduct{\nabla_X^2f(X)[H]}{H}_F.
    \]
    Nonnegative curvature in every direction is the second-order meaning of
    convexity.

    \item For $f(X)=\frac12\norm{AXB-C}_F^2$,
    \[
        \nabla_X^2f(X)[H]=A^\top AHB B^\top,
        \qquad
        \innerproduct{\nabla_X^2f(X)[H]}{H}_F=\norm{AHB}_F^2\ge0.
    \]
    Hence least-squares objectives are convex, and solving $\nabla_Xf=0$ gives
    a global minimizer in the unconstrained case.
\end{enumerate}
\end{gbox}






\subsection{Applications}

\subsubsection*{Deriving Estimators}

Consider the linear regression model written in matrix form
\[
    \mathbf{y}=X\boldsymbol{\beta}+\boldsymbol{\varepsilon},
\]
where $\mathbf{y}\in\mathbb{R}^n$, $X\in M_{n\times k}(\mathbb{R})$, $\boldsymbol{\beta}\in\mathbb{R}^k$, and $\boldsymbol{\varepsilon}\in\mathbb{R}^n$.
The sum of squared errors is
\[
    SSE(\boldsymbol{\beta})
    =
    \sum_{i=1}^n(y_i-\mathbf{x}_i^{\top}\boldsymbol{\beta})^2
    =
    \norm{\mathbf{y}-X\boldsymbol{\beta}}^2.
\]
Here $\mathbf{x}_i^{\top}$ is the $i$-th row of $X$.

\begin{theorem}[OLS estimator]
If $\rank(X)=k$, then
\[
    \widehat{\boldsymbol{\beta}}_{OLS}=(X^{\top}X)^{-1}X^{\top}\mathbf{y}.
\]
\end{theorem}

\begin{remark}
    This is just a special case of proposition~\eqref{affine least square}.
\end{remark}

\begin{pfbox}
Write
\[
    SSE(\boldsymbol{\beta})= \norm{X\boldsymbol{\beta}-\mathbf{y}}^2 = \tr((X\boldsymbol{\beta}-\mathbf{y})^{\top}(X\boldsymbol{\beta}-\mathbf{y})).
\]
Since $d(X\boldsymbol{\beta})=X\,d\boldsymbol{\beta}$,
\[
\begin{aligned}
    dSSE(\boldsymbol{\beta})
    &=
    \tr(d(X\boldsymbol{\beta})^{\top}(X\boldsymbol{\beta}-\mathbf{y}))
    +
    \tr((X\boldsymbol{\beta}-\mathbf{y})^{\top}d(X\boldsymbol{\beta}))\\
    &=
    \tr((X\,d\boldsymbol{\beta})^{\top}(X\boldsymbol{\beta}-\mathbf{y}))
    +
    \tr((X\boldsymbol{\beta}-\mathbf{y})^{\top}X\,d\boldsymbol{\beta})\\
    &=
    2\tr((X\boldsymbol{\beta}-\mathbf{y})^{\top}X\,d\boldsymbol{\beta}).
\end{aligned}
\]
Therefore
\[
    \nabla_{\boldsymbol{\beta}} SSE
    =
    \bigl(2(X\boldsymbol{\beta}-\mathbf{y})^{\top}X\bigr)^{\top}
    =
    2X^{\top}(X\boldsymbol{\beta}-\mathbf{y}).
\]
At the minimizer,
\[
    2X^{\top}(X\widehat{\boldsymbol{\beta}}-\mathbf{y})=0.
\]
Thus
\[
    X^{\top}X\widehat{\boldsymbol{\beta}}=X^{\top}\mathbf{y}.
\]
Since $\rank(X)=k$, the matrix $X^{\top}X$ is invertible. Hence
\[
    \widehat{\boldsymbol{\beta}}_{OLS}=(X^{\top}X)^{-1}X^{\top}\mathbf{y}.
\]

The matrix equation is the same as the summation version. Indeed,
\[
    X^{\top}X=\sum_{i=1}^n \mathbf{x}_i\mathbf{x}_i^{\top},
    \qquad
    X^{\top}\mathbf{y}=\sum_{i=1}^n \mathbf{x}_i y_i.
\]
Therefore,\\
the normal equation $X^{\top}X\widehat{\boldsymbol{\beta}}=X^{\top}\mathbf{y}$ is exactly $\left(\displaystyle\sum_{i=1}^n \mathbf{x}_i\mathbf{x}_i^{\top}\right)\widehat{\boldsymbol{\beta}}=\displaystyle\sum_{i=1}^n \mathbf{x}_i y_i$.
\end{pfbox}

\begin{remark}
The notation $SSE(\boldsymbol{\beta})$ emphasizes the squared-error criterion as
a function of an arbitrary coefficient vector:
\[
    SSE(\boldsymbol{\beta})
    =
    \sum_{i=1}^n
    (y_i-\mathbf{x}_i^{\top}\boldsymbol{\beta})^2.
\]
After an estimator $\widehat{\boldsymbol{\beta}}$ is chosen, the same quantity
is often called the residual sum of squares:
\[
    RSS
    =
    RSS(\widehat{\boldsymbol{\beta}})
    =
    SSE(\widehat{\boldsymbol{\beta}})
    =
    \sum_{i=1}^n \widehat e_i^2.
\]
Thus $SSE$ emphasizes the objective function, while $RSS$ emphasizes the fitted
residuals. Many texts use the two names interchangeably.
\end{remark}





\begin{definition}[Likelihood and log-likelihood]
For fixed data $z_1,\ldots,z_n$ with density or probability mass function
$f(z;\theta)$, define
\[
    L(\theta)=\prod_{i=1}^n f(z_i;\theta),
    \qquad
    \mathcal{L}(\theta)=\log L(\theta)
    =
    \sum_{i=1}^n \log f(z_i;\theta).
\]
Since $\log$ is increasing, maximizing $L(\theta)$ is equivalent to maximizing
$\mathcal{L}(\theta)$.
\end{definition}


\begin{theorem}[Normal Regression]
Condition on fixed regressors $\mathbf{x}_1,\ldots,\mathbf{x}_n\in\mathbb{R}^k$.
Suppose
\[
    y_i=\mathbf{x}_i^{\top}\boldsymbol{\beta}+e_i,
    \qquad
    e_i\iid N(0,\sigma^2).
\]
Equivalently, $y_i\iid\mathbf{x}_i\sim N(\mathbf{x}_i^{\top}\boldsymbol{\beta},\sigma^2)$. If $\sum_{i=1}^n\mathbf{x}_i\mathbf{x}_i^{\top}$ is invertible, then
\[
    \widehat{\boldsymbol{\beta}}_{MLE}=\widehat{\boldsymbol{\beta}}_{OLS}
    =
    \left(\sum_{i=1}^n \mathbf{x}_i\mathbf{x}_i^{\top}\right)^{-1}
    \left(\sum_{i=1}^n \mathbf{x}_i y_i\right),
    \qquad
    \widehat\sigma^2_{MLE}
    =
    \frac{1}{n}\sum_{i=1}^n\widehat e_i^2,
\]
where
\[
    \widehat e_i=y_i-\mathbf{x}_i^{\top}\widehat{\boldsymbol{\beta}}_{MLE}.
\]
\end{theorem}

\begin{pfbox}
For one observation,
\[
    f(y_i\mid\mathbf{x}_i;\boldsymbol{\beta},\sigma^2)
    =
    \frac{1}{\sqrt{2\pi\sigma^2}}
    \exp\left(
        -\frac{1}{2\sigma^2}
        (y_i-\mathbf{x}_i^{\top}\boldsymbol{\beta})^2
    \right).
\]
Thus
\[
    L(\boldsymbol{\beta},\sigma^2)
    =
    \prod_{i=1}^n f(y_i\mid\mathbf{x}_i;\boldsymbol{\beta},\sigma^2),
\]
and the log-likelihood is, up to an additive constant $C$,
\[
    \mathcal{L}(\boldsymbol{\beta},\sigma^2)
    =
    C
    -
    \frac{n}{2}\log\sigma^2
    -
    \frac{1}{2\sigma^2}
    \sum_{i=1}^n(y_i-\mathbf{x}_i^{\top}\boldsymbol{\beta})^2.
\]

For fixed $\sigma^2$, put $r_i=\mathbf{x}_i^{\top}\boldsymbol{\beta}-y_i$.
Then $(y_i-\mathbf{x}_i^{\top}\boldsymbol{\beta})^2=r_i^2$ and
$dr_i=\mathbf{x}_i^{\top}d\boldsymbol{\beta}$. Since $r_i$ is a scalar,
\[
\begin{aligned}
    d(r_i^2)
    &=
    d\,\tr(r_i^{\top}r_i)                                      \\
    &=
    \tr((dr_i)^{\top}r_i)+\tr(r_i^{\top}dr_i)                    \\
    &=
    2r_i\mathbf{x}_i^{\top}d\boldsymbol{\beta}.
\end{aligned}
\]
Therefore
\[
\begin{aligned}
    d_{\boldsymbol{\beta}}\mathcal{L}
    &=
    -\frac{1}{2\sigma^2}\sum_{i=1}^n
    2r_i\mathbf{x}_i^{\top}d\boldsymbol{\beta}                  \\
    &=
    \tr\left(
        -\frac{1}{\sigma^2}
        \sum_{i=1}^n r_i\mathbf{x}_i^{\top}
        d\boldsymbol{\beta}
    \right).
\end{aligned}
\]
Comparing with
$d_{\boldsymbol{\beta}}\mathcal{L}=(\nabla_{\boldsymbol{\beta}}\mathcal{L})^{\top}d\boldsymbol{\beta}$,
we read off
\[
    (\nabla_{\boldsymbol{\beta}}\mathcal{L})^{\top}
    =
    -\frac{1}{\sigma^2}
    \sum_{i=1}^n r_i\mathbf{x}_i^{\top}.
\]
Equivalently,
\[
    \nabla_{\boldsymbol{\beta}}\mathcal{L}
    =
    -\frac{1}{\sigma^2}
    \sum_{i=1}^n \mathbf{x}_i r_i
    =
    -\frac{1}{\sigma^2}
    \sum_{i=1}^n \mathbf{x}_i
    (\mathbf{x}_i^{\top}\boldsymbol{\beta}-y_i).
\]
Hence the first-order condition is
\[
    -\frac{1}{\sigma^2}
    \sum_{i=1}^n\mathbf{x}_i(\mathbf{x}_i^{\top}\boldsymbol{\beta}-y_i)
    =
    0.
\]
Equivalently,
\[
    \left(
        \underbrace{
            \sum_{i=1}^n \mathbf{x}_i\mathbf{x}_i^{\top}
        }_{X^{\top}X}
    \right)\boldsymbol{\beta}
    =
    \underbrace{
        \sum_{i=1}^n \mathbf{x}_i y_i
    }_{X^{\top}\mathbf{y}}.
\]
Hence
\[
    \widehat{\boldsymbol{\beta}}_{MLE}
    =
    \left(\sum_{i=1}^n \mathbf{x}_i\mathbf{x}_i^{\top}\right)^{-1}
    \left(\sum_{i=1}^n \mathbf{x}_i y_i\right).
\]
This is exactly the OLS estimator.

Next fix $\boldsymbol{\beta}$ and put
\[
    RSS(\boldsymbol{\beta})=\sum_{i=1}^n(y_i-\mathbf{x}_i^{\top}\boldsymbol{\beta})^2.
\]
Ignoring constants,
\[
    \mathcal{L}(\boldsymbol{\beta},\sigma^2)
    =
    -\frac{n}{2}\log\sigma^2
    -
    \frac{1}{2\sigma^2}RSS(\boldsymbol{\beta}).
\]
Treating $\sigma^2$ as the scalar variable,
\[
    \frac{\partial \mathcal{L}}{\partial \sigma^2}
    =
    -\frac{n}{2}\frac{1}{\sigma^2}
    +
    \frac{1}{2}\frac{1}{(\sigma^2)^2}RSS(\boldsymbol{\beta}).
\]
The first-order condition is
\[
    -\frac{n}{2}\frac{1}{\sigma^2}
    +
    \frac{1}{2}\frac{1}{(\sigma^2)^2}RSS(\boldsymbol{\beta})
    =
    0.
\]
Multiplying by $2(\sigma^2)^2$ gives
\[
    -n\sigma^2+RSS(\boldsymbol{\beta})=0.
\]
Thus, at $\widehat{\boldsymbol{\beta}}_{MLE}$,
\[
    \widehat\sigma^2_{MLE}
    =
    \frac{1}{n}RSS(\widehat{\boldsymbol{\beta}}_{MLE})
    =
    \frac{1}{n}\sum_{i=1}^n\widehat e_i^2.
\]
\end{pfbox}



\begin{theorem}[MLE Estimator]
Let $\mathbf{x}_1,\ldots,\mathbf{x}_n\iid N(\boldsymbol{\mu},\Sigma)$, where
$\mathbf{x}_i\in\mathbb{R}^p$, $\boldsymbol{\mu}\in\mathbb{R}^p$, and
$\Sigma\in M_p(\mathbb{R})$ is symmetric positive definite. Put
\[
    \bar{\mathbf{x}}
    =
    \frac{1}{n}\sum_{i=1}^n\mathbf{x}_i,
    \qquad
    S
    =
    \sum_{i=1}^n
    (\mathbf{x}_i-\bar{\mathbf{x}})
    (\mathbf{x}_i-\bar{\mathbf{x}})^{\top}.
\]
If $S$ is positive definite, then
\[
    \widehat{\boldsymbol{\mu}}_{MLE}
    =
    \bar{\mathbf{x}},
    \qquad
    \widehat\Sigma_{MLE}
    =
    \frac{1}{n}S.
\]
\end{theorem}

\begin{pfbox}
For one observation,
\[
    p(\mathbf{x}_i;\boldsymbol{\mu},\Sigma)
    =
    \frac{1}{(2\pi)^{p/2}(\det\Sigma)^{1/2}}
    \exp\left(
        -\frac{1}{2}
        (\mathbf{x}_i-\boldsymbol{\mu})^{\top}
        \Sigma^{-1}
        (\mathbf{x}_i-\boldsymbol{\mu})
    \right).
\]
Hence
\[
    L(\boldsymbol{\mu},\Sigma)
    =
    \prod_{i=1}^n p(\mathbf{x}_i;\boldsymbol{\mu},\Sigma),
\]
and the log-likelihood is
\[
    \mathcal{L}(\boldsymbol{\mu},\Sigma)
    =
    C
    -
    \frac{n}{2}\log\det\Sigma
    -
    \frac{1}{2}\sum_{i=1}^n
    (\mathbf{x}_i-\boldsymbol{\mu})^{\top}
    \Sigma^{-1}
    (\mathbf{x}_i-\boldsymbol{\mu}).
\]
Here $C$ does not depend on $\boldsymbol{\mu}$ or $\Sigma$.

First fix $\Sigma$. Since $\Sigma^{-1}$ is symmetric,
\[
\begin{aligned}
    d\left[
        (\mathbf{x}_i-\boldsymbol{\mu})^{\top}
        \Sigma^{-1}
        (\mathbf{x}_i-\boldsymbol{\mu})
    \right]
    &=
    -(d\boldsymbol{\mu})^{\top}\Sigma^{-1}(\mathbf{x}_i-\boldsymbol{\mu})
    -
    (\mathbf{x}_i-\boldsymbol{\mu})^{\top}\Sigma^{-1}d\boldsymbol{\mu} \\
    &=
    -2(\mathbf{x}_i-\boldsymbol{\mu})^{\top}\Sigma^{-1}d\boldsymbol{\mu}.
\end{aligned}
\]
Therefore
\[
\begin{aligned}
    d_{\boldsymbol{\mu}}\mathcal{L}
    &=
    \sum_{i=1}^n
    (\mathbf{x}_i-\boldsymbol{\mu})^{\top}\Sigma^{-1}d\boldsymbol{\mu} \\
    &=
    \left(
        \Sigma^{-1}\sum_{i=1}^n(\mathbf{x}_i-\boldsymbol{\mu})
    \right)^{\top}d\boldsymbol{\mu}.
\end{aligned}
\]
Comparing with
$d_{\boldsymbol{\mu}}\mathcal{L}=(\nabla_{\boldsymbol{\mu}}\mathcal{L})^{\top}d\boldsymbol{\mu}$,
we get
\[
    \nabla_{\boldsymbol{\mu}}\mathcal{L}
    =
    \Sigma^{-1}\sum_{i=1}^n(\mathbf{x}_i-\boldsymbol{\mu}).
\]
Since $\Sigma$ is invertible, the first-order condition gives
\[
    \sum_{i=1}^n\mathbf{x}_i-n\boldsymbol{\mu}=0,
    \qquad
    \widehat{\boldsymbol{\mu}}_{MLE}
    =
    \frac{1}{n}\sum_{i=1}^n \mathbf{x}_i.
\]

Now plug in $\widehat{\boldsymbol{\mu}}_{MLE}=\bar{\mathbf{x}}$.
For each $i$,
\[
\begin{aligned}
    &(\mathbf{x}_i-\widehat{\boldsymbol{\mu}}_{MLE})^{\top}
    \Sigma^{-1}
    (\mathbf{x}_i-\widehat{\boldsymbol{\mu}}_{MLE})                 \\
    &\qquad =
    \tr\left[
        (\mathbf{x}_i-\widehat{\boldsymbol{\mu}}_{MLE})^{\top}
        \Sigma^{-1}
        (\mathbf{x}_i-\widehat{\boldsymbol{\mu}}_{MLE})
    \right]                                                        \\
    &\qquad =
    \tr\left[
        \Sigma^{-1}
        (\mathbf{x}_i-\widehat{\boldsymbol{\mu}}_{MLE})
        (\mathbf{x}_i-\widehat{\boldsymbol{\mu}}_{MLE})^{\top}
    \right].
\end{aligned}
\]
Hence, with
\[
    S
    =
    \sum_{i=1}^n
    (\mathbf{x}_i-\widehat{\boldsymbol{\mu}}_{MLE})
    (\mathbf{x}_i-\widehat{\boldsymbol{\mu}}_{MLE})^{\top},
\]
we have
\[
\begin{aligned}
    \sum_{i=1}^n
    d\left[
        (\mathbf{x}_i-\widehat{\boldsymbol{\mu}}_{MLE})^{\top}
        \Sigma^{-1}
        (\mathbf{x}_i-\widehat{\boldsymbol{\mu}}_{MLE})
    \right]
    &=
    d\,\tr(\Sigma^{-1}S)                                      \\
    &=
    \tr\bigl(d(\Sigma^{-1})S\bigr)                             \\
    &=
    \tr\bigl(-\Sigma^{-1}(d\Sigma)\Sigma^{-1}S\bigr)            \\
    &=
    -\tr\bigl(\Sigma^{-1}S\Sigma^{-1}d\Sigma\bigr).
\end{aligned}
\]
where we used $d(\Sigma^{-1})=-\Sigma^{-1}(d\Sigma)\Sigma^{-1}$.
Together with $d(\log\det\Sigma)=\tr(\Sigma^{-1}d\Sigma)$, this gives
\[
    d\mathcal{L}
    =
    \tr\left[
        \left(
            -\frac{n}{2}\Sigma^{-1}
            +
            \frac{1}{2}\Sigma^{-1}S\Sigma^{-1}
        \right)d\Sigma
    \right].
\]
Hence the first-order condition is
\[
    -\frac{n}{2}\Sigma^{-1}
    +
    \frac{1}{2}\Sigma^{-1}S\Sigma^{-1}=0.
\]
Multiplying by $2$, then by $\Sigma$ on the left and by $\Sigma$ on the right,
gives
\[
    -n\Sigma+S=0.
\]
Therefore
\[
    \widehat\Sigma_{MLE}
    =
    \frac{1}{n}S
    =
    \frac{1}{n}\sum_{i=1}^n
    (\mathbf{x}_i-\widehat{\boldsymbol{\mu}}_{MLE})
    (\mathbf{x}_i-\widehat{\boldsymbol{\mu}}_{MLE})^{\top}.
\]
\end{pfbox}

\begin{remark}
Let $\mathbf{x}_1,\ldots,\mathbf{x}_n\in\mathbb{R}^p$, and let
\[
    S
    =
    \sum_{i=1}^n
    (\mathbf{x}_i-\bar{\mathbf{x}})
    (\mathbf{x}_i-\bar{\mathbf{x}})^\top
\]
be the centered scatter matrix. The covariance MLE is
\[
    \widehat{\Sigma}
    =
    \frac{1}{n}S.
\]
If $S$ is singular, then $\widehat{\Sigma}$ is also singular, so it is not
positive definite. Hence there is no covariance MLE inside the positive
definite cone.

In particular, since
\[
    \sum_{i=1}^n(\mathbf{x}_i-\bar{\mathbf{x}})=0,
\]
the centered vectors $\mathbf{x}_i-\bar{\mathbf{x}}$ span a space of dimension
at most $n-1$. Therefore
\[
    \rank(S)\le n-1.
\]
If $n\le p$, then $n-1<p$, so $S$ cannot have full rank $p$. Thus $S$ cannot
be positive definite. Hence a necessary condition for a positive definite
covariance MLE in $\mathbb{R}^p$ is $n>p$: the sample size must be larger than
the dimension, i.e. larger than the number of coordinates to be estimated
together.
\end{remark}









\subsubsection*{Low-rank approximation}

\begin{question}
Let $A\in M_{m\times n}(\mathbb{R})$ and $r\le\rank(A)$. Find $B\in M_{m\times n}(\mathbb{R})$ with $\rank(B)=r$ such that $\norm{A-B}_F$ is minimized.
\end{question}

A rank-$r$ matrix can be written in the form
\[
    B=XZ^{\top},
    \qquad
    X\in M_{m\times r}(\mathbb{R}),
    \qquad
    Z\in M_{n\times r}(\mathbb{R}).
\]
To see the meaning of this factorization, write $Z=(\mathbf{z}_1,\ldots,\mathbf{z}_r)$ with $\mathbf{z}_j\in\mathbb{R}^n$.
Then
\[
    Z^{\top}=
    \begin{pmatrix}
        \mathbf{z}_1^{\top}\\
        \vdots\\
        \mathbf{z}_r^{\top}
    \end{pmatrix}.
\]
If $X_i$ denotes the $i$-th row of $X$, then the $i$-th row of $B$ is $B_{i\cdot}=X_iZ^{\top}$,
which is a linear combination of the rows $\mathbf{z}_1^{\top},\ldots,\mathbf{z}_r^{\top}$.
Therefore, the row space of $B$ is spanned by $\{\mathbf{z}_1^{\top},\ldots,\mathbf{z}_r^{\top}\}$.

There is redundancy in this description. If $Q\in M_r(\mathbb{R})$ is invertible, then
\[
    XZ^{\top}=(XQ)(ZQ^{-\top})^{\top}.
\]
So $X$ and $Z$ are not uniquely determined by $B$. To remove part of this redundancy, impose the constraint
\[
    Z^{\top}Z=I_r.
\]
This means the columns $\mathbf{z}_1,\ldots,\mathbf{z}_r$ of $Z$ are orthonormal. Notice that $Z$ itself is not an orthogonal matrix unless $r=n$, because $Z$ is generally not square.

Thus we study the constrained optimization problem
\[
    \text{minimize } f(X,Z)=\frac{1}{2}\norm{A-XZ^{\top}}_F^2
    \qquad\text{subject to }Z^{\top}Z=I_r.
\]
The factor $1/2$ is inserted only to simplify derivatives later. 

Writing $Z=(\mathbf{z}_1,\ldots,\mathbf{z}_r)$, the constraint $Z^{\top}Z=I_r$ means
\[
    \mathbf{z}_i^{\top}\mathbf{z}_j=\delta_{ij}
    \qquad
    \text{for all }1\le i,j\le r.
\]
The scalar constraints are the entries of $Z^{\top}Z-I_r$. We organize the
corresponding Lagrange multipliers into an $r\times r$ matrix
$\Lambda=(\lambda_{ij})$. The trace term in the Lagrangian is the Frobenius
inner product form of all multiplier terms at once. Indeed, if
$C=Z^{\top}Z-I_r$, then $\tr(\Lambda C)=\innerproduct{\Lambda^{\top}}{C}_F$.
Since $C$ is symmetric, we may choose $\Lambda$ symmetric. Then
\[
    \tr\bigl(\Lambda(Z^{\top}Z-I_r)\bigr)
    =
    \innerproduct{\Lambda}{Z^{\top}Z-I_r}_F
    =
    \sum_{i=1}^r\sum_{j=1}^r
    \lambda_{ij}
    (\mathbf{z}_i^{\top}\mathbf{z}_j-\delta_{ij}).
\]
Thus $\frac{1}{2}\tr\bigl(\Lambda(Z^{\top}Z-I_r)\bigr)$ packages all scalar
constraints $\mathbf{z}_i^{\top}\mathbf{z}_j=\delta_{ij}$ and their Lagrange
multipliers into one matrix expression. Define
\[
    \mathcal{L}(X,Z,\Lambda)
    =
    \frac{1}{2}\tr(RR^{\top})
    -
    \frac{1}{2}\tr\bigl(\Lambda(Z^{\top}Z-I_r)\bigr),
\]
where $R=A-XZ^{\top}$.

\paragraph*{Step 1: First order condition in $X$}

\begin{theorem}[First order condition]
When $\mathcal{L}$ is minimized, one has $X=AZ$.
\end{theorem}

\begin{pfbox}
Regard $Z$ and $\Lambda$ as constants. Since the constraint term does not
depend on $X$ and $R=A-XZ^{\top}$, we have $d_XR=-dX\,Z^{\top}$. By the
chain rule and the differential of the squared Frobenius norm,
\[
\begin{aligned}
    d_X\mathcal{L}
    &=
    \frac12d_X\tr(RR^{\top})\\
    &=
    d_X\left(\frac12\norm{R}_F^2\right)\\
    &=
    d_R\left(\frac12\norm{R}_F^2\right)[d_XR]\\
    &=
    \innerproduct{R}{-dX\,Z^{\top}}_F\\
    &=
    \tr\left(R^{\top}(-dX\,Z^{\top})\right)\\
    &=
    -\tr(Z^{\top}R^{\top}dX)
    =
    -\tr((RZ)^{\top}dX).
\end{aligned}
\]
Hence $\nabla_X\mathcal{L}=-RZ$. At a minimizer,
$0=RZ=(A-XZ^{\top})Z=AZ-X$, where $Z^{\top}Z=I_r$ is used in the last
equality. Thus $X=AZ$.
\end{pfbox}

\paragraph*{Step 2: Reduction to a problem in $Z$}

\begin{theorem}
When $X=AZ$, the objective function becomes
\[
    f(AZ,Z)
    =
    \frac{1}{2}\left(\tr(A^{\top}A)-\tr(Z^{\top}A^{\top}AZ)\right).
\]
Therefore minimizing $f(AZ,Z)$ subject to $Z^{\top}Z=I_r$ is equivalent to maximizing
\[
    \tr(Z^{\top}A^{\top}AZ)
    \qquad\text{subject to }Z^{\top}Z=I_r.
\]
\end{theorem}

\begin{pfbox}
At a minimizer in $X$, we have $X=AZ$, so
$XZ^{\top}=AZZ^{\top}$. Put $P=ZZ^{\top}$. Since $Z^{\top}Z=I_r$,
$P^{\top}=P$ and $P^2=Z(Z^{\top}Z)Z^{\top}=ZZ^{\top}=P$.

Now $f(AZ,Z)=\frac12\norm{A-AZZ^{\top}}_F^2
=\frac12\norm{A(I-P)}_F^2$. Using $P^{\top}=P$ and $(I-P)^2=I-P$,
\[
\begin{aligned}
    f(AZ,Z)
    &=
    \frac{1}{2}\tr\bigl((A(I-P))^{\top}A(I-P)\bigr)\\
    &=
    \frac{1}{2}\tr\bigl((I-P)A^{\top}A(I-P)\bigr)\\
    &=
    \frac{1}{2}\tr\bigl(A^{\top}A(I-P)^2\bigr)\\
    &=
    \frac{1}{2}\tr\bigl(A^{\top}A(I-P)\bigr)\\
    &=
    \frac{1}{2}\left(\tr(A^{\top}A)-\tr(A^{\top}AP)\right).
\end{aligned}
\]
Since $P=ZZ^{\top}$, cyclicity of trace gives
$\tr(A^{\top}AP)=\tr(Z^{\top}A^{\top}AZ)$. Therefore
\[
    f(AZ,Z)
    =
    \frac{1}{2}\left(\tr(A^{\top}A)-\tr(Z^{\top}A^{\top}AZ)\right).
\]
The first term $\tr(A^{\top}A)$ does not depend on $Z$. Hence minimizing $f(AZ,Z)$ is equivalent to maximizing $\tr(Z^{\top}A^{\top}AZ)$ subject to $Z^{\top}Z=I_r$.
\end{pfbox}



\begin{theorem}
Let $P\in M_n(\mathbb{R})$. If $P^\top=P$ and $P^2=P$, then $P$ is the
orthogonal projection onto $\Ima(P)$.
\end{theorem}

\begin{pfbox}
Suppose $P^\top=P$ and $P^2=P$.

First, $P$ acts as the identity on $\Ima(P)$. Indeed, if
$\mathbf{y}\in\Ima(P)$, then $\mathbf{y}=P\mathbf{x}$ for some $\mathbf{x}$.
Hence $P\mathbf{y}=P^2\mathbf{x}=P\mathbf{x}=\mathbf{y}$.

Also, $P$ kills every vector in $\Ker(P)$. If $\mathbf{y}\in\Ker(P)$, then
$P\mathbf{y}=0$ by definition.

It remains to show that these two subspaces are orthogonal. Let
$\mathbf{u}\in\Ima(P)$ and $\mathbf{v}\in\Ker(P)$. Then
$\mathbf{u}=P\mathbf{x}$ for some $\mathbf{x}$, so
\[
\begin{aligned}
    \langle \mathbf{u},\mathbf{v}\rangle=\langle P\mathbf{x},\mathbf{v}\rangle=\langle \mathbf{x},P^\top\mathbf{v}\rangle=\langle \mathbf{x},P\mathbf{v}\rangle=0.
\end{aligned}
\]
Thus $\Ima(P)\perp\Ker(P)$.

Since $P^2=P$, every $\mathbf{x}\in\mathbb{R}^n$ can be decomposed as
$\mathbf{x}=P\mathbf{x}+(\mathbf{x}-P\mathbf{x})$, where
$P\mathbf{x}\in\Ima(P)$ and
\[
    P(\mathbf{x}-P\mathbf{x})
    =
    P\mathbf{x}-P^2\mathbf{x}
    =
    0.
\]
Hence $\mathbf{x}-P\mathbf{x}\in\Ker(P)$, and therefore
$\mathbb{R}^n=\Ima(P)\oplus\Ker(P)$. By orthogonal decomposition, we have
$\Ker(P)=\Ima(P)^\perp$.

Therefore $P$ preserves $\Ima(P)$ and kills $\Ima(P)^\perp$. Hence $P$ is the
orthogonal projection onto $\Ima(P)$.
\end{pfbox}

\begin{remark}
Idempotent alone is not enough to guarantee orthogonal projection. For example,
let
$$
P=
\begin{pmatrix}
    1 & 1\\
    0 & 0
\end{pmatrix}.
$$
Then $P^2=P$, so $P$ is idempotent. However, $P^\top\neq P$, so it is not
symmetric.

Indeed, $\Ima(P)=\operatorname{span}\{(1,0)^t\}$ and
$\Ker(P)=\operatorname{span}\{(1,-1)^t\}$. These two subspaces are not
orthogonal, since $\langle (1,0)^t,(1,-1)^t\rangle=1$. Hence $P$ is a
projection, but not an orthogonal projection.
\end{remark}

\begin{remark}
In our case, $P=ZZ^\top$. Since $Z^\top Z=I_r$, we have
$P^\top=P$ and
$$
P^2
=
ZZ^\top ZZ^\top
=
Z(Z^\top Z)Z^\top
=
ZZ^\top
=
P.
$$
Thus $P$ is symmetric and idempotent. By the theorem above, $P$ is the
orthogonal projection onto $\Ima(P)$.

Moreover, $\Ima(P)$ is exactly the column space of $Z$. Indeed, from
$P=ZZ^\top$, it is clear that $\Ima(P)$ is contained in the column space of
$Z$. Conversely, if $\mathbf{y}$ belongs to the column space of $Z$, then
$\mathbf{y}=Z\mathbf{a}$ for some $\mathbf{a}$, and
$$
P\mathbf{y}
=
ZZ^\top Z\mathbf{a}
=
Z\mathbf{a}
=
\mathbf{y}.
$$
Thus $\mathbf{y}\in\Ima(P)$. Therefore $\Ima(P)$ is the column space of $Z$.

So $P=ZZ^\top$ is the orthogonal projection onto the column space of $Z$.
\end{remark}






\paragraph*{Step 3: Linear algebra input}

\begin{theorem}[Rayleigh]
Let $B$ be symmetric and positive semidefinite, with eigenvalues
$\lambda_1\ge\lambda_2\ge\cdots\ge\lambda_n\ge 0$.
Then the maximum value of $\tr(Z^{\top}BZ)$ subject to $Z^{\top}Z=I_r$ is $\lambda_1+\cdots+\lambda_r$.

The maximum is attained when the columns of $Z$ are eigenvectors corresponding to the $r$ largest eigenvalues of $B$.
\end{theorem}

\begin{pfbox}
By the spectral theorem, there is an orthonormal basis $\mathbf{v}_1,\ldots,\mathbf{v}_n$ of eigenvectors of $B$ such that
$B\mathbf{v}_i=\lambda_i\mathbf{v}_i$. Write the columns of $Z$ as
$\mathbf{z}_1,\ldots,\mathbf{z}_r$. Since
$\mathbf{v}_1,\ldots,\mathbf{v}_n$ is an orthonormal basis, write
$\mathbf{z}_j=\displaystyle\sum_{i=1}^n\alpha_{ij}\mathbf{v}_i$. Then
\[
\begin{aligned}
    \tr(Z^{\top}BZ)
    &=
    \sum_{j=1}^r \mathbf{z}_j^{\top}B\mathbf{z}_j\\
    &=
    \sum_{j=1}^r
    \left\langle
        B\left(\sum_{i=1}^n\alpha_{ij}\mathbf{v}_i\right),
        \sum_{k=1}^n\alpha_{kj}\mathbf{v}_k
    \right\rangle\\
    &=
    \sum_{j=1}^r
    \left\langle
        \sum_{i=1}^n\lambda_i\alpha_{ij}\mathbf{v}_i,
        \sum_{k=1}^n\alpha_{kj}\mathbf{v}_k
    \right\rangle\\
    &=
    \sum_{j=1}^r\sum_{i=1}^n\lambda_i\alpha_{ij}^2.
\end{aligned}
\]
Define $\gamma_i=\displaystyle\sum_{j=1}^r\alpha_{ij}^2$. Then
$\tr(Z^{\top}BZ)=\displaystyle\sum_{i=1}^n\lambda_i\gamma_i$.
We need two observations about $\gamma_i$.
\begin{itemize}
    \item Since the columns of $Z$ are orthonormal,
    \[
    \begin{aligned}
        \sum_{i=1}^n\gamma_i
        &=
        \sum_{i=1}^n\sum_{j=1}^r\alpha_{ij}^2
        =
        \sum_{j=1}^r\sum_{i=1}^n\alpha_{ij}^2\\
        &=
        \sum_{j=1}^r\norm{\mathbf{z}_j}^2
        =
        r.
    \end{aligned}
    \]

    \item Let
    $V_r=\spn\{\mathbf{z}_1,\ldots,\mathbf{z}_r\}$. Since
    $\mathbf{z}_1,\ldots,\mathbf{z}_r$ is an orthonormal basis of $V_r$,
    \[
        P_{V_r}(\mathbf{v}_i)
        =
        \sum_{j=1}^r
        \innerproduct{\mathbf{v}_i}{\mathbf{z}_j}\mathbf{z}_j.
    \]
    Therefore
    \[
    \begin{aligned}
        \norm{P_{V_r}(\mathbf{v}_i)}^2
        &=
        \sum_{j=1}^r
        \left|
            \innerproduct{\mathbf{v}_i}{\mathbf{z}_j}
        \right|^2\\
        &=
        \sum_{j=1}^r
        \left|
            \innerproduct{\mathbf{z}_j}{\mathbf{v}_i}
        \right|^2
        =
        \sum_{j=1}^r\alpha_{ij}^2
        =
        \gamma_i.
    \end{aligned}
    \]
    By the minimal-distance property of orthogonal projection,
    \[
        0
        \le
        \gamma_i
        =
        \norm{P_{V_r}(\mathbf{v}_i)}^2
        \le
        \norm{\mathbf{v}_i}^2
        =
        1
        \qquad
        \text{for every }1\le i\le n.
    \]
\end{itemize}

It remains to maximize
$\displaystyle\sum_{i=1}^n\lambda_i\gamma_i$ subject to
$0\le\gamma_i\le1$ and $\displaystyle\sum_{i=1}^n\gamma_i=r$.
Since $\lambda_1\ge\cdots\ge\lambda_n$, the maximum is achieved by placing all available weight on the largest eigenvalues:
$\gamma_1=\cdots=\gamma_r=1$ and
$\gamma_{r+1}=\cdots=\gamma_n=0$. Thus the maximum value is
$\lambda_1+\cdots+\lambda_r$, attained by taking
$Z=(\mathbf{v}_1,\ldots,\mathbf{v}_r)$.
\end{pfbox}

\begin{theorem}[Low-rank approximation theorem]
Let $A=U\Sigma V^{\top}$ be a singular value decomposition of $A$, with singular values
$\sigma_1\ge\sigma_2\ge\cdots\ge 0$. Then the best rank-$r$ approximation to $A$ in Frobenius norm is the truncated SVD
\[
    B_r
    =
    \sum_{i=1}^r \sigma_i \mathbf{u}_i\mathbf{v}_i^{\top}.
\]
Equivalently, if $U_r=(\mathbf{u}_1,\ldots,\mathbf{u}_r)$, $V_r=(\mathbf{v}_1,\ldots,\mathbf{v}_r)$, and $\Sigma_r=\operatorname{diag}(\sigma_1,\ldots,\sigma_r)$, then
\[
    B_r=U_r\Sigma_rV_r^{\top}.
\]
\end{theorem}

\begin{pfbox}
From Step 2, the problem reduces to maximizing
\[
    \tr(Z^{\top}A^{\top}AZ)
    \qquad\text{subject to }Z^{\top}Z=I_r.
\]
The matrix $A^{\top}A$ is symmetric and positive semidefinite. Let
\[
    A^{\top}A\mathbf{v}_i=\lambda_i \mathbf{v}_i,
    \qquad
    \lambda_1\ge\lambda_2\ge\cdots\ge\lambda_n\ge 0,
\]
where $\mathbf{v}_1,\ldots,\mathbf{v}_n$ is an orthonormal basis. The eigenvalues of $A^{\top}A$ are the squares of the singular values of $A$, so $\lambda_i=\sigma_i^2$.
By Rayleigh's theorem, the maximum of $\tr(Z^{\top}A^{\top}AZ)$ is
\[
    \lambda_1+\cdots+\lambda_r
    =
    \sigma_1^2+\cdots+\sigma_r^2,
\]
and it is attained when $Z=(\mathbf{v}_1,\ldots,\mathbf{v}_r)=V_r$.
From Step 1, the corresponding $X$ is $X=AZ=AV_r$.
Write $U=(U_r\ U_{\perp})$. Since the columns of $V$ are orthonormal,
\[
    V^{\top}V_r
    =
    \begin{pmatrix}
        I_r\\
        \mathbf O_{(n-r)\times r}
    \end{pmatrix}.
\]
Thus multiplication by $V^{\top}V_r$ selects the first $r$ columns of
$\Sigma$, and hence
\[
    \Sigma V^{\top}V_r
    =
    \begin{pmatrix}
        \Sigma_r\\
        \mathbf O_{(m-r)\times r}
    \end{pmatrix}.
\]
Therefore
\[
\begin{aligned}
    B
    &=
    XZ^{\top}
    =
    AV_rV_r^{\top}\\
    &=
    U\Sigma V^{\top}V_rV_r^{\top}\\
    &=
    (U_r\ U_{\perp})
    \begin{pmatrix}
        \Sigma_r\\
        \mathbf O_{(m-r)\times r}
    \end{pmatrix}
    V_r^{\top}\\
    &=
    U_r\Sigma_rV_r^{\top}
    =
    \sum_{i=1}^r \sigma_i \mathbf{u}_i\mathbf{v}_i^{\top}.
\end{aligned}
\]
Thus the best rank-$r$ approximation is the truncated SVD.
\end{pfbox}









